\chapter{安全证书、博弈-CBF 融合与多智能体强化学习基础}
\label{ch:plan-safemarl}

% ════════════════════════════════════════════════════════════════
%  全吸收章：融合两份 Tutorial 笔记
%   - 04_移动机器人规控/40_博弈规划/50_安全证书与MARL.md（CBF+博弈、GCBF+、势博弈、PSRO、CFR、OpenSpiel）
%   - 04_移动机器人规控/50_多机器人协作/110_MARL基础.md（Markov Game/Dec-POMDP、非平稳性、CTDE、值分解、策略梯度）
%  记号归一：R∈SO(3)、T∈SE(3)（preamble 宏 \SO \SE）；规划/控制/博弈自然记号保留。
%  跨部去重：CBF/CLF、HJ 可达性、鲁棒/随机 MPC 的最深推导留待【控制部】，本章侧重规划/决策用途。
%  源由 AI 生成多轮审校，存疑处就地 \pz / \rebuilt，并入 claimsToVerify。
%  教学层遵循 v5.0 理论教学规范。
% ════════════════════════════════════════════════════════════════

\rebuilt{本章融合自两份 Tutorial 笔记（博弈规划"安全证书与 MARL"、多机器人协作"MARL 基础"），其源稿由 AI 撰写并经多轮审校，但论文归属、年份、会议与实验数值仍可能有误。逐处 \texttt{\textbackslash cite} 标源；正文存疑处已就地 \texttt{\textbackslash pz}/\texttt{\textbackslash rebuilt} 标注，待联网核对后二次校验。}

\begin{practice}[前置自测]
开始之前，先试答下列问题。若已能流畅作答，可只速读本章表格与速查卡；若卡住，正是本章要为你解决的。
\begin{enumerate}
  \item 给定安全集 $\mathcal{C}=\{x:h(x)\ge0\}$，控制屏障函数（CBF）条件 $\dot h(x,u)+\alpha(h(x))\ge0$ 保证了什么？为什么对控制仿射系统它能写成一个关于 $u$ 的二次规划（QP）？
  \item 什么是纯策略 Nash 均衡、什么是混合策略 Nash 均衡？为什么石头剪刀布没有纯策略 Nash，却一定有混合策略 Nash？
  \item 写出 Q-learning 的更新式。值函数 $Q(s,a)$ 与策略 $\pi(a\mid s)$ 是什么关系？这个更新收敛到唯一不动点，依赖转移核 $P(s'\mid s,a)$ 的什么性质？
  \item 把单智能体强化学习算法（如 DQN/PPO）复制 $N$ 份、让每个机器人各自独立训练，会出什么问题？为什么"调小学习率"治不了它？
  \item 什么是 Double Oracle 算法？给定一个巨大的博弈，为什么"维护一个小策略子集、轮流加入对当前混合策略的最佳响应"能收敛到原博弈的 Nash？
\end{enumerate}
\noindent\emph{答案线索}：1$\to$\cref{sec:sm-cbf-game}（CBF-QP 与 GNE）；2$\to$\cref{sec:sm-marl-form}（Nash 与均衡分类）；3$\to$\cref{sec:sm-nonstationary}（非平稳性破坏压缩映射）；4$\to$\cref{sec:sm-nonstationary,sec:sm-ctde}；5$\to$\cref{sec:sm-psro}（PSRO 把 Double Oracle 推广到深度 RL）。
\end{practice}

\section*{本章目标}
学完本章，你应当能够：
\begin{enumerate}
  \item \textbf{说清}两种安全的分工——基于预测的"软安全"与基于不变集的"硬安全"，并理解控制屏障证书为何不依赖对手意图模型；
  \item \textbf{推导}多机最小范数 CBF-QP，并\textbf{论证}这组 QP 的 KKT 条件联合等价于一个广义 Nash 均衡（GNE）——"安全证书"自动成为"博弈均衡"的一部分；
  \item \textbf{解释} GCBF+ 如何用图神经网络同时学出图控制屏障函数与分布式控制器，以及"成对碰撞的局部性"为何让单一证书可覆盖任意规模；
  \item \textbf{掌握}势博弈与势函数的定义，\textbf{论证}"势函数存在 $\Rightarrow$ 纯 Nash 存在 $\Rightarrow$ 分布式学习收敛"，并说明 Voronoi 覆盖与 Markov 势博弈如何为合作 MARL 的收敛性奠基；
  \item \textbf{严格写出}合作 MARL 的两套形式化——随机博弈（Markov Game）与去中心化部分可观测 MDP（Dec-POMDP），并区分合作/竞争/混合三类；
  \item \textbf{从数学上论证} MARL 的第一性困难——非平稳性，说明它如何破坏压缩映射与经验回放，并理解 CTDE 如何"把漂移从环境挪进条件"；
  \item \textbf{从零推导}值分解线（VDN/QMIX）的 IGM 原则与单调混合，以及策略梯度线（MADDPG/COMA/MAPPO/HAPPO）的反事实基线与多智能体优势分解引理；
  \item \textbf{理解} MARL 与博弈论交界的三主线（值分解、Stackelberg RL、种群博弈），并\textbf{吃透} PSRO 如何用可插拔的 meta-solver 统一独立 RL、虚拟自博弈、Double Oracle、JPSRO；
  \item \textbf{会做技术选型}：面对一个多智能体问题，按"博弈结构 $\times$ 规模"两条轴判断该用 CBF+博弈、势博弈、MARL/PSRO 还是 OpenSpiel/CFR，并说清各自代价。
\end{enumerate}

\subsection*{本章知识导航}
本章的主线是一句话——\textbf{安全控制、多机协调、多智能体学习，本质都是博弈，只是博弈的类型不同}。把"博弈"当作贯穿全章的统一语言，四个工具就挂在同一棵树上：

\begin{center}
\small
\begin{tikzpicture}[
  node distance=5mm and 7mm, >=Stealth,
  blk/.style={draw, rounded corners, align=center, font=\small, inner sep=3pt, fill=main!6, draw=main!60!black},
  grp/.style={draw, rounded corners, align=center, font=\small, inner sep=3pt, fill=second!6, draw=second!60!black},
  adv/.style={draw, rounded corners, align=center, font=\small, inner sep=3pt, fill=third!8, draw=third!60!black},
  ln/.style={->, thick, gray!60!black}]
\node[blk] (root) {博弈作为统一语言\\ 安全控制 $\leftrightarrow$ 博弈 $\leftrightarrow$ 多智能体学习};
\node[grp, below left=10mm and -6mm of root] (cbf) {CBF + 博弈\\ CBF-QP 的 KKT $=$ GNE\\ GCBF+ 用 GNN 学证书};
\node[grp, below=10mm of root] (pot) {势博弈\\ 势函数 $\to$ 纯 NE\\ $\to$ 分布式收敛};
\node[grp, below right=10mm and -6mm of root] (marl) {MARL + PSRO\\ Markov Game/Dec-POMDP\\ CTDE / 值分解 / 策略梯度};
\node[blk, below=26mm of root, fill=main!8, draw=main!60!black] (sel) {统一框架与选型\\ OpenSpiel / CFR；离散 vs 连续\\ 结构 $\times$ 规模两条轴};
\draw[ln] (root)--(cbf); \draw[ln] (root)--(pot); \draw[ln] (root)--(marl);
\draw[ln] (cbf.south) |- (sel.west);
\draw[ln] (pot)--(sel);
\draw[ln] (marl.south) |- (sel.east);
\draw[ln, dashed] (pot.east) -- node[midway,font=\scriptsize,above]{MPG} (marl.west);
\end{tikzpicture}
\end{center}

\noindent\textbf{推荐路径}：\cref{sec:sm-cbf-game}（CBF + 博弈）是本章的"安全地基"，给出不依赖对手模型的硬安全证书，建议先读；\cref{sec:sm-potential}（势博弈）相对独立、偏理论，是多机协调的基石，且通过 Markov 势博弈连到合作 MARL 的收敛性；\cref{sec:sm-marl-form}（形式化）与 \cref{sec:sm-nonstationary}（非平稳性）、\cref{sec:sm-ctde}（CTDE）是 MARL 的公共语言，必读；\cref{sec:sm-valuedecomp}（值分解）与 \cref{sec:sm-policygrad}（策略梯度）是两条硬核技术路线，值得花最多时间；\cref{sec:sm-psro}（三主线与 PSRO）承接前面、把博弈与深度学习焊接；\cref{sec:sm-frameworks}（OpenSpiel/CFR 与选型）是收口的基础设施与设计空间。带（选读）色彩的 GCBF+ 理论细节、QTRAN/QPLEX、CFR 家族在首读时可略过、抓住主干思想。

\subsection*{前置知识桥接}
本章重度依赖控制部的安全/最优控制工具与规划部的搜索/动态规划，这里把会反复取用的几个结论激活一下，重叠处一律 \cref 指过去、不重复：

\begin{itemize}
  \item \textbf{控制屏障函数与 Lyapunov 稳定性}（控制部）：CBF 的思想脱胎于 Lyapunov 的"前向不变"——\cref{sec:ctrl-lyapunov} 给出了 Lyapunov 稳定/渐近稳定/指数稳定三层定义，\cref{thm:ctrl-lyap-direct} 是直接法。本章 \cref{sec:sm-cbf-game} 把"安全集前向不变"写成 CBF 不等式，势博弈的收敛性论证（\cref{sec:sm-potential}）也复用"势函数 $=$ Lyapunov 函数、单调下降必收敛"这一逻辑。\pz{控制部深讲：CBF/CLF 的完整理论（高相对度、输入约束、relaxed-CBF 可行性）与 HJ 可达性、鲁棒/随机 MPC 的安全滤波统一视角，应在控制部另设章深讲；本章只给规划/决策所需的关键式与直觉，待补 \cref。}
  \item \textbf{Bellman 最优性与动态规划}（\cref{sec:plan-dp}，\cref{thm:plan-bellman}）：MARL 把单步决策"序贯化"，其值函数与压缩映射收敛性直接建在 Bellman 递推之上；本章 \cref{sec:sm-nonstationary} 论证的核心，正是非平稳性如何破坏 \cref{thm:plan-bellman} 收敛所依赖的"转移核固定"前提。
  \item \textbf{LQR/MPC 与最优控制}（\cref{sec:ctrl-lqr,sec:ctrl-mpc}）：连续微分博弈的迭代求解（iLQR 风格）与 MPC 的滚动时域同源；本章把它与离散扩展式博弈（OpenSpiel）对照（\cref{sec:sm-frameworks}）。
  \item \textbf{搜索与采样式规划}（\cref{sec:plan-search,sec:plan-sampling}）：本章的多机安全证书是规划器之上的"安全滤波层"——规划器（搜索/采样/优化，见 \cref{ch:planning}）给标称控制，CBF 层做最小修正保安全。
\end{itemize}

\begin{note}
\textbf{本章记号与约定.}\;
旋转矩阵记 $\mathbf{R}\in\SO(3)$、位姿 $\mathbf{T}\in\SE(3)$（preamble 宏 $\SO,\SE,\so,\se$，与全书一致，\nameref{ch:notation}）；多机/控制/博弈领域的自然记号保留。
智能体数 $N$，下标 $i,j$，$-i$ 表"除 $i$ 外的所有智能体"；联合动作 $\mathbf{a}=(a_1,\dots,a_N)$、联合策略 $\boldsymbol{\pi}=(\pi_1,\dots,\pi_N)$、分布式策略 $\pi_i(a_i\mid o_i)$ 或 $\pi_i(a_i\mid\tau_i)$（$\tau_i$ 为局部动作-观测历史）。
机器人 $i$ 的状态 $x_i$、控制 $u_i$、位置 $p_i$；控制仿射动力学 $\dot x_i=f_i(x_i)+g_i(x_i)u_i$。
安全集 $\mathcal{C}=\{x:h(x)\ge0\}$，CBF 标量函数 $h$，扩展 class-$\mathcal{K}$ 函数 $\alpha$（线性化时记系数 $\alpha>0$）。
\textbf{字母隔离声明}：本章 $\alpha$ 指 CBF 的 class-$\mathcal{K}$ 系数（非学习率）、$\Phi$ 指势函数（非位姿）、$\sigma$ 指经验博弈上的混合 meta 策略（非奇异值或协方差）、$Q$ 在 MARL 语境指动作值函数、$V$ 指状态值函数，均仅本章局部生效。学习率记 $\eta$，折扣因子 $\gamma$。
\end{note}

% ════════════════════════════════════════════════════════════════
\section{CBF + 博弈：多智能体安全证书}
\label{sec:sm-cbf-game}

\paragraph{动机：推断可能错，但碰撞不能发生。}
设想一个交互式规划器（如自动驾驶博弈规划）：你在线推断对手的代价、用"预测即均衡"规划一条期望最优且安全的轨迹\pz{博弈规划部分（逆博弈、iLQGames、Level-k、预测即均衡）应在规划部另设博弈规划章 G1–G3 深讲；本章承接其"安全兜底"需求，待这些章写就后改 \cref 指过去。}。绝大多数时候它表现很好，但隐患始终悬着——\textbf{均衡的正确性依赖对手模型正确，而推断可能错}：对手突然加速抢行、而你的模型以为它会让，均衡预测就错了，规划出的动作可能把你和它送上碰撞航线。当推断可能错、而错误会致命时，光有"基于模型的最优"不够。

这里的关键区分是\textbf{两种安全}。"预测即均衡"给的是\emph{软安全}：基于对对手的最优预测规划，安全性和预测准确性绑死——预测错，安全也无从谈起。但有些后果（碰撞、撞人）\textbf{绝对不能发生}，不能赌在"预测足够准"上。我们需要一层\emph{硬安全}：无论对手实际怎么动、无论预测对不对，都保证不碰撞。

控制屏障函数（CBF）恰好提供这种硬安全。它不预测对手要干什么，而是给出一个\textbf{不变集}——只要系统状态留在安全集内、且控制满足 CBF 约束，状态就\emph{永远}不会离开安全集（前向不变）。对碰撞避免，安全集就是"机间距大于安全距离"的集合；CBF 约束保证一旦机器人之间足够远，它们就永远不会靠得太近。这个保证是\emph{几何}的（关于不变集），不依赖任何对手意图模型——正是它能给推断错误兜底的原因。

\paragraph{反面：为什么不能只靠"预测 + 避让"。}
一个看似够用的替代方案：用预测器预测对手轨迹，规划时把预测轨迹当障碍绕开。它有两个根本问题。其一，\textbf{安全性和预测准确性绑死}——你绕开的是\emph{预测的}对手轨迹，不是\emph{实际的}；对手没按预测走（抢行了），你绕开的是一条它没走的轨迹，反而可能撞上它实际走的。其二，\textbf{没有形式化保证}——"把预测当障碍绕开"是启发式，大多数情况好用但\emph{无法证明}永不碰撞，不能用在安全攸关的系统里。

CBF 的不同在于：只要 CBF-QP 在每一步可行（有满足约束的控制），就\emph{可证明}安全集前向不变、永不碰撞——这是一个定理，不是经验。代价是它\textbf{只管安全、不管最优}：CBF 只保证不碰撞，不保证高效或舒适，那是标称控制器（规划器）的事。所以正确的架构是\textbf{两层}：规划器给标称控制（追求最优），CBF 层在其上做最小修正（保证安全）——下面会看到这正是 CBF-QP 的形式。

\paragraph{历史：从单机 CBF-QP 到多机安全证书。}
CBF 的核心思想可追溯到 Ames 等人把控制 Lyapunov 函数的"前向不变"思路用于安全\cite{ames2017cbf}：用标量函数 $h(x)$ 刻画安全集 $\mathcal{C}=\{x:h(x)\ge0\}$，并要求控制让 $h$ 在边界上不减太快。把它推广到\textbf{多机器人碰撞避免}的奠基工作，是 Wang、Ames、Egerstedt 的"无碰撞多机系统的安全屏障证书"\cite{wang2017safety}。

\begin{paper}[Wang--Ames--Egerstedt 2017：多机安全屏障证书]\label{paper:wang2017}
\textbf{问题}：多机器人系统如何在保持标称任务的同时形式化地保证可证明无碰撞、且可扩展？
\textbf{核心思想}：对标称控制器做\emph{最小侵入式}（minimally invasive）修正以满足安全约束。三件关键事：(1) 把碰撞避免写成 CBF 约束——任意两机器人定义 $h_{ij}=\|p_i-p_j\|^2-D_s^2$（机间距平方减安全距离平方），$h_{ij}\ge0$ 即安全；(2) 把"尽量不打扰原计划"写成最小范数目标 $\min\|u-u_{\text{nom}}\|^2$（最小侵入）；(3) 证明这个 QP 可\emph{去中心化}——每个机器人只需对\emph{邻近}机器人保持安全、用局部信息解一个小 QP。
\textbf{关键结果}：保守性、解的存在性、死锁分别用 relaxed CBF、混合制动控制器、一致扰动（consistent perturbations）处理。
\textbf{与本书关系}：本节的多机 CBF-QP（\cref{eq:sm-cbf-qp}）即出自此，并据此建立 CBF-QP 与广义 Nash 均衡的等价（\cref{thm:sm-cbf-gne}）。
\textbf{局限}：死锁在纯 CBF 框架里无法根除，是"局部安全保证"换来的代价（\cref{sec:sm-cbf-game} 末讨论）。
\end{paper}

这条线后来枝繁叶茂：Notomista--Egerstedt 把它推广为"约束驱动协调"（用 CBF 约束编码长期自主任务）\cite{notomista2019constraint}；而把整套证书用\textbf{图神经网络学出来、扩展到上千 agent} 的，是本节要精读的 GCBF+\cite{zhang2025gcbfplus}。

\begin{insight}[CBF $\leftrightarrow$ 可达性 $\leftrightarrow$ 不变集]
CBF 给出的"安全集前向不变"，和 Hamilton--Jacobi（HJ）可达性是\emph{同一件事的两面}。相似：两者都在刻画"系统能否永远留在安全集"——HJ 算"避免集的补"（哪些状态无论如何都会进入危险），CBF 给"安全集前向不变的充分条件"。不同：HJ 要解一个偏微分方程（贵、精确、全局），CBF 只需在当前状态解一个 QP（快、保守、局部）。一句话：\textbf{HJ 是"精确但算不动"的安全验证，CBF 是"保守但实时"的安全综合}。在多机场景，HJ 的维度随 agent 数爆炸完全不可行，而 CBF-QP 的去中心化让它能上车——这是 CBF 成为多机安全主流的根本原因。\pz{控制部深讲：HJ 可达性偏微分方程、值函数即 CBF（PNCBF）、CBF/HJ/MPC 安全滤波统一，应在控制部安全控制章给完整推导，待 \cref。}
\end{insight}

\subsection{理论：CBF-QP 的 KKT 条件就是一个广义 Nash 均衡}
\label{sec:sm-cbf-gne}

现在到本节最漂亮的洞察——\textbf{多机 CBF-QP 与广义 Nash 均衡（GNE）的等价}。先写出单机 CBF-QP。机器人 $i$ 的状态 $x_i$、控制 $u_i$，控制仿射动力学 $\dot x_i=f_i(x_i)+g_i(x_i)u_i$。对每个邻居 $j$，安全约束 $h_{ij}(x_i,x_j)\ge0$ 的 CBF 条件（沿轨迹安全裕度不减太快）是
\begin{equation}
  \dot h_{ij}=\nabla_{x_i}h_{ij}\cdot\dot x_i+\nabla_{x_j}h_{ij}\cdot\dot x_j\ \ge\ -\alpha\,h_{ij},
  \label{eq:sm-cbf-cond}
\end{equation}
其中 $\alpha>0$ 是扩展 class-$\mathcal{K}$ 函数的（线性）系数\cite{wang2017safety}。机器人 $i$ 在尽量贴近标称控制 $u_{i,\text{nom}}$（来自规划器）的前提下满足所有邻居的安全约束：
\begin{equation}
  u_i^*=\arg\min_{u_i}\ \|u_i-u_{i,\text{nom}}\|^2\quad\text{s.t.}\quad
  \nabla_{x_i}h_{ij}\cdot\big(f_i+g_iu_i\big)+\alpha\,h_{ij}\ \ge\ -\nabla_{x_j}h_{ij}\cdot\dot x_j,\ \ \forall j\in\mathcal{N}_i.
  \label{eq:sm-cbf-qp}
\end{equation}
注意约束右边含 $\dot x_j$——即\textbf{邻居 $j$ 的运动}。这就是耦合的来源：机器人 $i$ 的可行控制集，依赖邻居 $j$ 怎么动；反过来 $j$ 的可行集也依赖 $i$。对控制仿射系统，约束左边关于 $u_i$ 是\emph{线性不等式}，故 \cref{eq:sm-cbf-qp} 是凸二次规划（CBF-QP）。

\begin{insight}[共享约束下的并行优化 $=$ 广义 Nash 均衡]
盯住 \cref{eq:sm-cbf-qp} 的耦合结构。$N$ 个机器人，每个解一个优化问题（min 自己的目标），但每个的\emph{约束依赖他人的决策}（$h_{ij}$ 同时含 $x_i,x_j$ 和 $\dot x_j$）。这正是\textbf{广义 Nash 均衡问题（GNEP）}的定义：多个玩家各自优化，玩家的\emph{可行域}（而不只是目标）依赖他人策略。普通 Nash：玩家目标耦合、可行域独立；广义 Nash：连可行域都耦合（通过共享约束）。CBF 的安全约束 $h_{ij}$ 是机器人 $i,j$ \emph{共享}的——它同时约束两者，正是 GNEP 的"共享约束"。于是\textbf{多机 CBF-QP 的解（每个机器人的最优安全控制）$=$ 这个 GNEP 的一个广义 Nash 均衡}。
\end{insight}

\begin{theorem}[多机 CBF-QP 的 KKT 等价于 GNE 的 KKT]\label{thm:sm-cbf-gne}
设 $N$ 个机器人各解 \cref{eq:sm-cbf-qp}，乘子记 $\lambda_{ij}\ge0$。把每个 QP 的 Karush--Kuhn--Tucker（KKT）条件——稳定性（拉格朗日量对 $u_i$ 梯度为零）、原始可行（满足 CBF 约束）、对偶可行（$\lambda_{ij}\ge0$）、互补松弛（$\lambda_{ij}\,[\text{约束裕度}]_{ij}=0$）——按 $i$ 堆叠，所得方程组的未知量恰为全体 $\{u_i\}$ 与全体 $\{\lambda_{ij}\}$，方程数与之匹配，构成一个混合互补问题（MCP），其结构与把 GNEP 写成 MCP 的统一方程组完全一致。因此这组 CBF-QP 的联合解是该 GNEP 的一个广义 Nash 均衡。
\end{theorem}
\begin{proof}[证明思路]\rebuilt{（据 \cite{wang2017safety} 与广义 Nash/MCP 标准结果）}
单个 QP \cref{eq:sm-cbf-qp} 是凸的，强对偶成立，其最优解由 KKT 条件\emph{充要}刻画。机器人 $i$ 的拉格朗日量
$L_i=\|u_i-u_{i,\text{nom}}\|^2-\sum_{j}\lambda_{ij}\big(\nabla_{x_i}h_{ij}\cdot(f_i+g_iu_i)+\alpha h_{ij}+\nabla_{x_j}h_{ij}\cdot\dot x_j\big)$，
稳定性条件 $\partial L_i/\partial u_i=2(u_i-u_{i,\text{nom}})-\sum_j\lambda_{ij}\,g_i^\top\nabla_{x_i}h_{ij}=0$。约束里出现 $\dot x_j=f_j+g_ju_j$——即\emph{他人决策} $u_j$，这正是耦合项。把 $N$ 组（稳定性 $+$ 原始/对偶可行 $+$ 互补松弛）叠在一起，得到的方程组里每个 $u_i$ 的方程含 $\{u_j\}_{j\in\mathcal{N}_i}$、每个互补松弛含共享 $h_{ij}$——这恰是 GNEP 经 KKT 化为 MCP 的标准形式 $F(z)=0,\ 0\le\lambda\perp G(z)\ge0$。故联合 KKT 解即一个 GNE。
\end{proof}

这个等价不是文字游戏。它的深层意义是——\textbf{当每个机器人都"自私地"求解自己的最小侵入安全 QP 时，它们集体达到的状态恰好是一个博弈均衡}。没有谁在做全局协调（去中心化），但共享的安全约束让它们的并行决策自动构成一个一致的均衡。这就是为什么称 CBF-QP 的解为"安全\emph{证书}"——它不仅是一个控制，还是一个均衡，自带"无人能单方面更安全地偏离"的博弈论性质。实践含义：你不需要为多机安全单独设计一个协调器；只要每个机器人解自己的 CBF-QP，安全就以博弈均衡的形式涌现。

\begin{note}
\textbf{理论 vs 实现.}\;
\cref{thm:sm-cbf-gne} 说"解\emph{是}一个 GNE"是\emph{理论性质}；实现上几乎总是\emph{去中心化近似}——每个机器人独立解自己的 QP，把邻居控制 $u_j$（即 $\dot x_j$）当\emph{已知}（用上一步或标称值代入）。这是一种 Jacobi/Gauss--Seidel 式并行近似，不是联立精确求解整个 GNE。两者不矛盾：去中心化迭代在温和条件下收敛到 GNE。误以为要把所有 QP 联立成一个大问题求解，会失去去中心化的可扩展优势。
\end{note}

\subsection{GCBF+：用图神经网络学出"任意规模的单一证书"}
\label{sec:sm-gcbfplus}

手工设计的 CBF-QP 有两个扩展瓶颈：其一，每个机器人要对\emph{所有}邻居写约束，邻居多时 QP 约束爆炸；其二，CBF 本身（$h_{ij}$ 的形式、$\alpha$ 的选择）要人工设计，复杂动力学（如四旋翼）下很难手工调好。把整套证书"学"出来、扩展到上千 agent 的核心论文是 GCBF+\cite{zhang2025gcbfplus}。

\begin{paper}[GCBF+：图控制屏障函数与分布式安全控制]\label{paper:gcbfplus}
\textbf{问题}：带障碍的大规模环境中，大量 agent 如何仅用\emph{局部信息}保持安全并到达目标？
\textbf{核心思想}：引入一类新证书——\emph{图控制屏障函数（GCBF）}，把整个多智能体系统建成一张\emph{图}（节点是机器人和障碍，边是"足够近、需关注"的邻居关系），每个机器人只看图上自己的\emph{局部邻域}（固定大小，不随总数增长）；用图神经网络（GNN）同时参数化候选 GCBF $h$（安全证书）与分布式控制策略 $\pi$，联合训练；可直接处理 LiDAR 点云而非真实状态。
\textbf{关键结果}：(1) 一个理论框架，用\emph{单一} GCBF 证明\emph{任意规模} MAS 的安全；(2) 8 个 agent 训练即可部署到 1024 个 agent；(3) Crazyflie 无人机群硬件实验。
\textbf{与本书关系}：它把 \cref{sec:sm-cbf-gne} 手工 CBF-QP 的"安全 $=$ 博弈均衡"原理，借学习扩展到工业规模——原理与规模的两端。
\textbf{局限}：邻域大小 $l$ 是敏感超参；连续时间 CBF 在离散步进下有数值 gap；理论安全建立在"控制不变集"的有限时间可达近似上。
\end{paper}

\paragraph{逐步讲解：局部安全如何聚合成任意规模全局安全。}
GCBF+ 最硬核的贡献是"用单一 GCBF 证明任意数量 agent 安全"。这听起来近乎魔法，拆开看逻辑很清晰，分三步：
\begin{enumerate}
  \item \textbf{定义"局部安全"}。传统 CBF 的安全集是全局的（所有 agent 的联合状态），维度随 agent 数爆炸。GCBF 换视角：定义每个 agent 在其\emph{局部邻域}（图上邻居 $+$ 邻近障碍）内的安全——agent $i$ 安全，当且仅当它与邻域内每个对象的距离都大于安全距离。这个局部安全的"维度"固定（邻域大小固定）。
  \item \textbf{局部安全 $\Rightarrow$ 全局安全}。关键观察：碰撞总是\emph{成对}发生的（agent $i$ 撞 agent $j$）。若每个 agent 都在自己邻域内安全（与所有邻居保持距离），则任何一对邻近 agent 都不碰撞；而非邻近 agent 本来就离得远（不在彼此邻域内 $=$ 距离大于邻域半径 $>$ 安全距离）。于是"每个 agent 局部安全"$\Rightarrow$"所有 agent 对都不碰撞"$\Rightarrow$"系统全局安全"。
  \item \textbf{单一 GCBF 编码局部安全}。既然全局安全可由局部安全聚合，那么只需学一个刻画"局部邻域安全"的 GCBF $h$——其输入是 agent 的局部邻域（固定维度的图），输出是局部安全裕度。同一个 $h$ 应用到每个 agent 的局部邻域，就覆盖了任意数量的 agent。这就是"单一证书、任意规模"的来源：\textbf{证书作用在固定大小的局部邻域上，而非随规模变化的全局状态上}。
\end{enumerate}

\begin{insight}["成对碰撞"是可聚合性的物理根源]
盯住第二步——为什么局部安全能聚合成全局安全？根本原因是\textbf{碰撞的"局部性"}：碰撞是两个物体的接触，是一个\emph{成对}事件，不存在"三个 agent 同时碰撞但两两都不碰"的情况。这个物理事实让安全可以分解：全局无碰撞 $\Leftrightarrow$ 所有 agent 对无碰撞 $\Leftrightarrow$ 每个 agent 与其邻居无碰撞（非邻居自动满足）。对比\emph{不可}这样分解的性质：比如"全局最优覆盖"就不能由局部聚合（一个 agent 的最优位置依赖全局其他所有 agent），所以覆盖问题（\cref{sec:sm-potential}）要用势博弈的全局 $\Phi$，而安全问题可以用局部 GCBF。\textbf{可分解的性质才能局部聚合，安全恰好可分解。}
\end{insight}

\paragraph{训练损失：把 CBF 的数学条件变成可学的目标。}
GCBF+ 联合训练两个 GNN（一个输出 GCBF 值 $h$、一个输出控制 $u$），关键是把 CBF 的数学条件翻译成神经网络损失，三部分各编码 CBF 的一个要求\cite{zhang2025gcbfplus}：
\begin{itemize}
  \item \textbf{安全集分类损失}：CBF 必须在安全状态为正、危险状态为负（$h(x)\ge0\Leftrightarrow x$ 安全）。采样安全状态（间距大）与危险状态（间距小或已碰撞），用类似分类的损失逼 $h$ 在前者为正、后者为负，使 $h$ 的零水平集对齐真实安全边界。
  \item \textbf{CBF 条件损失（导数条件）}：CBF 的核心是 $\dot h+\alpha h\ge0$。在每个采样状态上用当前控制器算 $\dot h$，用 hinge 损失惩罚违反者，使 $h$ 真正成为"控制下不变"的证书，而非随便一个分类器。
  \item \textbf{目标到达损失}：控制器不仅要安全，还要完成任务（到目标）。用到目标距离类损失，保证学到的策略不是"原地不动以保安全"（呼应死锁问题）。
\end{itemize}
三部分加权求和联合训练 $h$ 和 $u$——这是"安全证书 $+$ 控制器联合学习"的精髓：证书定义什么是安全，控制器学会在安全内完成任务，两者互相塑造。

\begin{note}
\textbf{关键实验数据（数量级，具体以论文为准）.}\;\rebuilt{以下数字转述自 \cite{zhang2025gcbfplus}，待联网核对。}
GCBF+ 在五个环境（2D：单积分器/双积分器/独轮车；3D：线性无人机/Crazyflie 四旋翼）验证；核心"规模实验"在无障碍环境固定邻域大小、把 agent 数 $N$ 从 8 增到 1024（密度增 100 倍以上）。报告称：中小规模（$\le256$ agents）安全率最高、手工 CBF-QP 约低 20\%；大规模（1024 agents）比领先 RL（如 MAPPO）约高 40\%；训练只用 8 个 agent、约 1000 步。\textbf{真正卖点是"不牺牲任务的安全"}——手工 CBF 把 $\alpha$ 调小也很安全但过保守（绕大圈、到不了），RL 难两全；GCBF+ 用"以 QP 控制器作 reference $+$ 三部分损失"打破"安全 vs 任务"权衡。看安全方法时永远要同时看"安全率"和"任务完成率"。\pz{存疑：256 agents 高 20\%、1024 agents 高 40\%、训练 1h vs MAPPO 5h 等具体数值需对照论文核实；GCBF（前身）vs GCBF+ 的数字勿张冠李戴。}
\end{note}

\begin{note}
\textbf{神经 CBF 的三条近期路线（对照）.}\;\rebuilt{年份/会议待核：}GCBF+\cite{zhang2025gcbfplus}（图 CBF $+$ GNN，攻\emph{可扩展性}，吃 LiDAR 点云）；PNCBF\cite{so2024pncbf}（学\emph{标称策略的值函数}作 CBF，攻\emph{高相对度 $+$ 输入约束下难手工构造}）；BarrierNet\cite{xiao2023barriernet}（把 CBF-QP 做成\emph{可微层}嵌入神经控制器，攻\emph{环境自适应、端到端可学}）。共同母题是"用学习突破手工 CBF 的设计瓶颈"，不同路线攻不同痛点——大规模 swarm 选 GCBF+，单体高维 $+$ 输入约束选 PNCBF，要端到端可学的安全层选 BarrierNet。
\end{note}

\subsection{代码走读：多机 CBF-QP 避碰（安全证书的最小实现）}
\label{sec:sm-cbf-code}

把前面的理论落成能跑的代码：4 个 2D 单积分器机器人做位置交换（position swap），必经中心交叉，CBF-QP 保证永不碰撞。为聚焦"安全证书"逻辑，动力学用最简单的单积分器 $\dot x_i=u_i$，真实工程里换成 GCBF+ 的学习式策略或更复杂动力学。每个机器人解一个最小范数 QP：在满足所有邻居 CBF 约束的前提下，控制尽量贴近标称（朝目标）。

\begin{codebox}[Python]{多机 CBF-QP 安全证书（最小实现）}
import numpy as np
from scipy.optimize import minimize
N, Ds, alpha, dt, umax = 4, 0.5, 3.0, 0.05, 1.5    # n_robots / safe_dist / CBF_gain / dt / u_bound

def nominal(xi, gi):                                # nominal control: head to goal (saturated)
    d = gi - xi; n = np.linalg.norm(d)
    return d / n * min(1.0, n) if n > 1e-6 else np.zeros(2)

def cbf_qp(i, X, u_nom_all):
    # robot i CBF-QP: min ||u - u_nom||^2 s.t. CBF constraint for each neighbor
    xi = X[i]; u_nom = u_nom_all[i]; cons = []
    for j in range(N):
        if j == i:
            continue
        e = xi - X[j]                               # relative position
        h = e @ e - Ds ** 2                         # safety margin h_ij = ||p_i - p_j||^2 - Ds^2
        # CBF: 2 e . (u_i - u_j) + alpha . h >= 0  ->  (2 e) . u_i >= -alpha h + 2 e . u_j
        A = 2 * e; b = -alpha * h + 2 * e @ u_nom_all[j]   # treat u_j as known (decentralized)
        cons.append({'type': 'ineq', 'fun': (lambda u, A=A, b=b: A @ u - b)})
    return minimize(lambda u: np.sum((u - u_nom) ** 2), u_nom, method='SLSQP',
                    constraints=cons, bounds=[(-umax, umax)] * 2,
                    options={'maxiter': 150, 'ftol': 1e-10}).x
\end{codebox}

仿真主循环把 4 个机器人放在正方形四角、目标设在各自对角（必经中心冲突）。每个机器人对每个邻居约束依赖邻居的运动（$h_{ij}$ 含 $u_j$），所以 $N$ 个投影相互耦合——它们的联合解就是 \cref{thm:sm-cbf-gne} 说的 GNE。"安全投影"是单机视角，"广义 Nash 均衡"是系统视角，两者是同一组 QP 的两种读法。

\paragraph{死锁：CBF 多机安全绕不开的实际问题。}
在对称场景（如四角对穿）下，纯 CBF-QP 会\textbf{死锁}——机器人完全安全（不碰撞）但僵在原地永远到不了目标。常见解死锁手段：(1) 对称破缺——加切向偏置/随机扰动或右手通行约定；(2) 优先级——给机器人排序，低优先级让高优先级；(3) 高层协调——用一个规划器分配通过顺序（牺牲一点去中心化）。GCBF+ 用学习式策略一定程度缓解死锁（学到的策略隐含绕行行为），但死锁在纯 CBF 框架里无法根本消除——这是"局部安全保证"换来的代价，也呼应 \cref{sec:sm-potential} 势博弈为何要引入"全局势函数"来保证收敛（活性）。

\paragraph{relaxed CBF：当硬约束无可行解时。}
密集场景下多个 CBF 约束相互冲突（尤其控制有界 $|u|\le u_{\max}$ 时），QP 可能无可行解，求解器返回失败/NaN——而这恰是最危险的时刻。修复用 \textbf{relaxed CBF}：把硬约束改成软约束 $+$ 松弛变量惩罚
\begin{equation}
  \min_{u,\,\delta}\ \|u-u_{\text{nom}}\|^2+M\sum_j\delta_{ij}^2
  \quad\text{s.t.}\quad \text{CBF}_{ij}\ge-\delta_{ij},\ \delta_{ij}\ge0,
  \label{eq:sm-relaxed-cbf}
\end{equation}
保证 QP 永远有解，松弛量 $\delta_{ij}$ 的大小指示安全裕度被侵犯的程度（当"安全裕度传感器"用），并配降级策略（松弛过大时紧急制动）。

\begin{pitfall}[概念误区：CBF 保证"一定到达目标"]
既然 CBF-QP 保证安全又追踪标称控制，新手常以为它既安全又一定能完成任务。实际上 CBF 保证的是\textbf{安全（safety，坏事不发生）}不是\textbf{活性（liveness，好事会发生）}——它只承诺"不碰撞"，从不承诺"到得了"。对称/密集场景会死锁，机器人完全安全但僵住，此时安全保证仍然成立，只是任务失败。正确做法：把 CBF 当\emph{安全滤波器}而非\emph{完整规划器}——到达目标是标称控制器（规划器）的职责，死锁要专门处理（对称破缺/优先级/高层协调）。检验方法：跑对称场景看是否死锁。
\end{pitfall}

\begin{pitfall}[思维陷阱：CBF 安全证书"不需要对手信息"]
CBF 是"不依赖对手\emph{意图}模型"的硬安全，于是有人以为它完全不需要任何关于邻居的信息。错。"不依赖对手意图/代价模型" $\neq$ "不依赖对手\emph{状态}信息"——约束 $h_{ij}$ 和 $\dot h_{ij}$ 都依赖邻居的位置和速度（\cref{eq:sm-cbf-qp} 右端含 $\dot x_j$）。CBF 免除的是\emph{意图建模}（不用逆博弈推断对手代价），但需要\emph{状态感知}（邻居的相对位置/速度，靠传感器或通信）。GCBF+ 正是用 LiDAR 点云提供这个状态信息。检验方法：列出 CBF 约束依赖的所有量，确认部署时可获得。
\end{pitfall}

\paragraph{软硬两层安全栈：预测 $+$ 证书如何配合。}
把本节放回完整交互系统，会看到一个两层安全架构，缺一不可：\emph{上层软安全}（基于预测的均衡规划，追求性能——舒适、高效、类人，但安全性建立在对手预测准确上）；\emph{下层硬安全}（CBF 安全证书，把上层轨迹当标称控制做最小修正，保证不变集前向不变，不依赖任何对手意图模型）。两层配合：绝大多数时候上层规划本就安全、下层不激活（输出 $=$ 标称）；只在上层因预测错误规划出危险动作时，下层才介入修正——把预测错误的后果从"碰撞"降级为"保守"。这是一个一般性的安全系统设计原则：\textbf{把基于模型的性能优化与不依赖模型的硬安全保证分层，用后者为前者的模型错误兜底}。

\begin{practice}
\begin{enumerate}
  \item 复现 4 机器人位置交换 demo，扩展到 8、16 个机器人，观察死锁是否更频繁、最小机间距如何变化；实现 relaxed CBF（\cref{eq:sm-relaxed-cbf}），对比硬约束版在密集场景的可行性。
  \item 写出 2 个机器人 CBF-QP 各自的 KKT 条件（含乘子、互补松弛），把它们堆叠，验证这个堆叠系统正是一个 GNE 的 KKT（对照 \cref{thm:sm-cbf-gne}）；解释共享约束 $h_{12}$ 如何让两个 QP 的 KKT 耦合。
  \item 扫描 CBF 系数 $\alpha\in\{0.5,1,3,10\}$：记录每个 $\alpha$ 下的最小机间距、是否到达目标、轨迹绕行程度，论证 $\alpha$ 小则保守（早减速、绕大圈、可能死锁）、$\alpha$ 大则激进（贴边界、险但高效），找一个平衡的 $\alpha$。
\end{enumerate}
\end{practice}

% ════════════════════════════════════════════════════════════════
\section{势博弈与多机器人：分布式协调何时收敛}
\label{sec:sm-potential}

\paragraph{动机：自私优化为什么通常不收敛。}
\cref{sec:sm-cbf-game} 的 CBF 保证安全但不保证活性（会死锁）。更一般地，多个 agent 各自优化时，集体行为何时会\textbf{收敛}到一个好的稳态？把多机协调建模成博弈，立刻撞上一个根本问题：博弈的均衡可能不存在（纯策略意义下）、或者分布式学习不收敛。回忆石头剪刀布——它没有纯策略 Nash，每个玩家的最佳响应循环往复（出石头$\to$对手出布$\to$我出剪刀$\to\dots$），永远转圈。一般博弈里，让每个 agent 反复对他人当前策略做最佳响应（最自然的分布式学习），完全可能陷入这种循环。我们需要一个\textbf{保证}：什么样的博弈，能确保自私的分布式优化一定收敛到一个稳定、且对全局有意义的均衡？答案是\textbf{势博弈}。

\paragraph{反面：一般博弈的最佳响应动力学会发散。}
设想你不知道势博弈，直接让多机做"最佳响应迭代"。一般博弈里这会出两种问题：其一，\emph{循环不收敛}——A 调整以适应 B，B 又调整以适应 A 的新位置，若博弈是"非势"的（像石头剪刀布的连续版），位置永远在变；其二，\emph{收敛了也不一定对全局好}——即便碰巧收敛到一个纯 Nash，它也只是"无人能单方面改进"的局部稳定点，多个 Nash 中可能有好有坏，自私优化可能卡在差的那个。势博弈同时解决这两点。

\paragraph{历史。}
势博弈的概念由 Monderer 和 Shapley 正式提出\cite{monderer1996potential}：在（精确）势博弈中，任一玩家单方面偏离所带来的\emph{自身}收益变化，恰好等于一个\emph{全局}势函数的变化。把它系统地用于\textbf{多机器人合作控制}的代表工作是 Marden、Arslan、Shamma\cite{marden2009cooperative}——用"博弈中的学习"语言看待合作控制，把一致性、动态传感器覆盖等问题表述成势博弈，于是分布式学习算法的收敛性由势博弈理论保证。

\subsection{理论：势函数、纯 Nash、分布式收敛}

\begin{definition}[精确势博弈]\label{def:sm-potential}
博弈 $(N,\{A_i\},\{J_i\})$ 是\textbf{精确势博弈}，若存在\textbf{势函数} $\Phi:A\to\mathbb{R}$，使对任意玩家 $i$、任意 $a_i,a_i'\in A_i$、任意他人策略 $a_{-i}$：
\begin{equation}
  J_i(a_i',a_{-i})-J_i(a_i,a_{-i})=\Phi(a_i',a_{-i})-\Phi(a_i,a_{-i}).
  \label{eq:sm-potential-def}
\end{equation}
（约定代价 $J_i$ 越小越好，玩家最小化 $J_i$ 等价于最小化 $\Phi$；若用收益则取最大化。）
\end{definition}

关键含义：每个玩家"自私地"改善自己的收益，\textbf{等价于}改善同一个全局函数 $\Phi$——就像每个人都在爬同一座山，虽然各爬各的，但都往山顶去。

\begin{theorem}[势博弈的三个核心性质]\label{thm:sm-potential-props}
设 $\Phi$ 是有限动作集上的势函数，则：
\begin{enumerate}
  \item \textbf{纯 Nash 存在性}：$\Phi$ 在有限动作集上必有最小值点 $a^*$；在 $a^*$ 处无人能单方面降低 $\Phi$（已是极小），由 \cref{eq:sm-potential-def} 即无人能单方面降低自己的 $J_i$，故 $a^*$ 是纯 Nash。势博弈\emph{一定}有纯策略 Nash 均衡（一般博弈不保证）。
  \item \textbf{更优响应收敛}：任何"单方面改善"（某玩家把 $J_i$ 降低）都让 $\Phi$ 下降同样多（由 \cref{eq:sm-potential-def}）；$\Phi$ 在有限集上有下界，故这样的改善序列必然终止于一个纯 Nash（$\Phi$ 不能无限下降）。分布式更优响应动力学\emph{一定收敛、不循环}。
  \item \textbf{虚拟博弈/梯度流收敛}：势博弈具有虚拟博弈（fictitious play）性质\cite{monderer1996potential}；连续动作下，势函数的梯度流（每个 agent 沿 $-\partial\Phi/\partial a_i$ 移动）收敛到 $\Phi$ 的临界点。
\end{enumerate}
\end{theorem}

\begin{insight}[势函数把"博弈不动点"变成"优化极值"]
盯住性质二——这是势博弈最实用的性质。一般博弈的最佳响应动力学可能循环（石头剪刀布），因为没有一个"单调下降的量"保证收敛。势博弈有 $\Phi$ 这个\textbf{Lyapunov 函数}：每次有人改善，$\Phi$ 严格下降；$\Phi$ 有下界，所以下降必然终止——这就排除了循环（与 \cref{sec:ctrl-lyapunov} 的 Lyapunov 收敛逻辑同源）。换句话说，$\Phi$ 给分布式学习提供了"势能"，系统像球滚下山一样必然停在谷底（纯 Nash $=$ $\Phi$ 局部极小）。这正是为什么多机协调要"设计成势博弈"：你不是在求一个可能不存在/不收敛的一般 Nash，而是在让多个 agent 集体优化一个有下界的 $\Phi$——收敛由 $\Phi$ 的存在性免费保证。
\end{insight}

\subsection{多机器人应用：Voronoi 覆盖}

最经典的势博弈式多机协调是\textbf{传感器覆盖}：$N$ 个机器人要分布在区域 $Q$ 上，使"每个点都离最近的机器人尽量近"。定义 locational cost（覆盖代价）作为势函数：
\begin{equation}
  \Phi(p_1,\dots,p_N)=\sum_{i=1}^N\int_{V_i}\|q-p_i\|^2\,\phi(q)\,\mathrm dq,
  \label{eq:sm-voronoi}
\end{equation}
其中 $V_i$ 是机器人 $i$ 的 Voronoi 格（所有离 $i$ 最近的点），$\phi(q)$ 是密度（哪里更重要）。这是一个势博弈：每个机器人移动只影响自己的 Voronoi 格那部分代价，单机改善 $=$ $\Phi$ 改善。最优策略是\textbf{移向自己 Voronoi 格的质心}（Lloyd 算法）——这正是 $\Phi$ 的梯度下降。

\begin{codebox}[Python]{Voronoi 覆盖作为势博弈（gradient play）}
import numpy as np
np.random.seed(0)
N = 5
gx, gy = np.meshgrid(np.linspace(0, 1, 60), np.linspace(0, 1, 60))
Q = np.column_stack([gx.ravel(), gy.ravel()])    # discretized region points (for integration)
phi = np.ones(len(Q))                            # importance density (uniform)
P = np.random.rand(N, 2) * 0.4 + 0.3             # init: 5 robots clustered at center

def assign(P):                                   # Voronoi partition: each point to nearest robot
    d = np.linalg.norm(Q[:, None, :] - P[None, :, :], axis=2)
    return np.argmin(d, axis=1)

def potential(P):                                # potential Phi: coverage cost
    lab = assign(P); J = 0.0
    for i in range(N):
        m = lab == i
        if m.any():
            J += np.sum(np.sum((Q[m] - P[i]) ** 2, axis=1) * phi[m])
    return J

for it in range(40):                             # gradient play = Lloyd iteration
    lab = assign(P); newP = P.copy()
    for i in range(N):
        m = lab == i
        if m.any():
            w = phi[m]
            newP[i] = np.sum(Q[m] * w[:, None], axis=0) / np.sum(w)   # move to Voronoi centroid
    P = newP                                     # Phi decreases monotonically -> pure Nash
\end{codebox}

跑出来势函数单调下降并收敛，所有机器人停在各自 Voronoi 格的质心——此时无人能单方面降低 $\Phi$（移动激励为零），系统到达纯 Nash。这就是势博弈威力的具体兑现：\textbf{5 个机器人各自自私地优化（移向自己格质心），集体却在单调降低同一个全局覆盖代价 $\Phi$，并保证收敛到纯 Nash}——没有中心协调器，收敛由势函数的存在性免费保证。

\begin{insight}[Lloyd 算法 $=$ 势函数的坐标下降]
盯住循环里"每个机器人移向 Voronoi 质心"这一步——它其实是势函数 $\Phi$ 的\textbf{坐标下降}（每次固定他人、优化自己那一维）。移向质心是单机最优响应（让自己负责的区域覆盖代价最小），由势博弈性质，这等价于降低全局 $\Phi$。所以经典 Lloyd 覆盖算法本质是势博弈上的分布式优化——这解释了它为什么总收敛（$\Phi$ 单调降 $+$ 有下界）、为什么去中心化能 work（每个机器人只需局部信息：自己的 Voronoi 格）。与 \cref{sec:sm-cbf-game} 的呼应：CBF-QP 是"安全 $=$ 博弈均衡"，这里是"协调 $=$ 博弈均衡的优化"——两者都把多机问题归结为博弈，只是前者求 GNE（安全约束耦合）、这里降势函数（合作目标对齐）。
\end{insight}

\paragraph{Voronoi 不止用于覆盖：Buffered Voronoi Cells 做避碰。}
Voronoi 划分还能直接用于多机\textbf{避碰}，连接覆盖与安全。代表工作是 Zhou、Wang、Bandyopadhyay、Schwager 的 Buffered Voronoi Cells\cite{zhou2017bvc}。核心观察：Voronoi 划分天然把空间分给各机器人、彼此不重叠，所以\emph{只要每个机器人待在自己的 Voronoi 格内，就不会和别人碰}。但机器人有物理尺寸，于是把 Voronoi 格边界\emph{向内缩进一个安全半径}，得到 \textbf{Buffered Voronoi Cell（BVC）}——只要机器人\emph{中心}在 BVC 内，它的\emph{整个身体}就在 Voronoi 格内，保证不碰。它只需感知相对位置（不需通信、不需速度），$O(k)$ 复杂度，适合带噪声传感器的在线部署。

\begin{insight}[BVC 与 CBF：两种"几何安全"的对照]
把 BVC 和 CBF 放一起，能看清两种"几何安全"思路的异同。相同：两者都给\emph{不依赖对手意图}的硬几何安全，都只需相对位置。不同：BVC 是"划分空间、各占一块"的\textbf{分割}思路（把安全变成"待在自己区域"）；CBF 是"约束控制、保持距离"的\textbf{约束}思路（把安全变成"满足不等式"）。BVC 对单积分器有简洁解析解；CBF 更一般（任意控制仿射系统、任意安全集）但需解 QP。它们和速度障碍（ORCA）一起构成不依赖意图建模的几何避碰家族，与学习式避碰互补：几何法有保证但偏保守，学习法更灵活但保证弱。
\end{insight}

\subsection{Markov 势博弈：合作 MARL 的收敛基石}
\label{sec:sm-mpg}

前面的势博弈都是\textbf{单步}（静态）博弈：agent 选一个动作、拿一次收益。但 MARL 里的博弈是\textbf{序贯}的——agent 在一个马尔可夫决策过程里反复选动作、状态不断转移。把势博弈从单步推广到序贯，就得到 \textbf{Markov 势博弈（Markov Potential Game, MPG）}，它是连接本节势博弈与 \cref{sec:sm-valuedecomp,sec:sm-policygrad} 合作 MARL 的关键理论桥梁。

\begin{definition}[Markov 势博弈]\label{def:sm-mpg}
随机博弈是 Markov 势博弈，若存在势函数 $\Phi(s,\boldsymbol{\pi})$（依赖状态 $s$ 和联合策略 $\boldsymbol{\pi}$），使任一 agent 单边改\emph{策略}时，它的\emph{价值函数}变化等于 $\Phi$ 的变化（即把 \cref{eq:sm-potential-def} 的"动作 $\to$ 策略、收益 $\to$ 价值函数"）。最重要的特例：\textbf{所有 agent 共享同一团队奖励}（合作 MARL）时，团队的价值函数本身就是势函数——故\emph{合作 MARL $=$ Markov 势博弈}（不是类比，是严格的数学事实）。
\end{definition}

\begin{insight}[MPG 为什么让"各学各的"也能收敛]
MPG 最深刻的推论是——\textbf{在 MPG 里，各 agent 独立做策略梯度（各学各的），也能收敛到 Nash}。这非常反直觉：\cref{sec:sm-nonstationary} 刚说过"独立 RL 因非平稳而不收敛"，为什么 MPG 是例外？因为 MPG 有那个势函数 $\Phi$ 作"全局 Lyapunov 函数"：每个 agent 的策略梯度改善自己的价值，由 MPG 性质等于改善 $\Phi$；所有 agent 一起爬同一个 $\Phi$，$\Phi$ 有界，故必然收敛——非平稳被 $\Phi$ 的单调性"驯服"了。近期理论把这点钉死：去中心化策略梯度在 MPG 全局收敛\cite{leonardos2022mpg}，且有不显式依赖状态空间大小的迭代复杂度\cite{ding2022mpg}。\textbf{合作 MARL 恰好落在这个"独立学习也能收敛"的安全区里（团队奖励 $=$ 势函数），这就是 VDN/QMIX/MAPPO 等合作 MARL 有收敛保证的根本原因。}
\end{insight}

\begin{note}
\textbf{势博弈划出的收敛分水岭.}\;
\emph{有势结构}一侧（合作 MARL、势博弈式协调）：独立/简单学习就收敛，有全局 $\Phi$ 兜底——这是本节和值分解主线的地盘。\emph{无势结构}一侧（一般和、竞争博弈）：独立 RL 不收敛（\cref{sec:sm-nonstationary}），必须用 PSRO（种群）、CFR（no-regret）等专门机制逼近均衡——这是 \cref{sec:sm-psro,sec:sm-frameworks} 的地盘。判断一个多智能体问题"好不好学"，第一刀就是看它有没有（近似）势结构。真实问题很少是\emph{精确}势博弈，近期工作放松到\emph{近似}势博弈（Markov $\alpha$-Potential Games），仍能保证收敛到近似 Nash（误差随 $\alpha$ 缩放）。
\end{note}

\begin{pitfall}[概念误区：任何多机协调问题都是势博弈]
看到势博弈这么好用（保证收敛），有人以为只要是合作的多机问题就天然是势博弈、自然就收敛。错。势博弈要求存在满足 \cref{eq:sm-potential-def} 的全局势函数——这是\emph{强条件}，不是所有博弈都满足。合作（共享目标）是充分条件之一，但很多多机问题（agent 目标不完全一致、有局部冲突）不是势博弈。正确做法：先\emph{验证}博弈是否有势函数（检查代价能否写成 $\Phi$ $+$ 与自身无关项），或在设计协调问题时\emph{主动}把目标设计成势函数（如用 locational cost）。另一个常见误区：收敛到纯 Nash 就以为是全局最优——纯 Nash $=$ $\Phi$ 的\emph{局部}极小，$\Phi$ 可能非凸有多个局部极小，需多次随机初始化/退火跳出差的局部极小（log-linear learning 用带温度的随机选择可收敛到全局最优）。
\end{pitfall}

\begin{practice}
\begin{enumerate}
  \item 复现 5 机器人 Voronoi 覆盖 demo：用非均匀密度 $\phi(q)$（某区域更重要），观察机器人是否向高密度区聚集；验证势函数单调下降；多次随机初始化，统计收敛势值的分布，讨论局部极小。
  \item 把多机任务分配（$N$ 机器人分配到 $M$ 任务点，最小化总距离）建模成势博弈：设计势函数（总分配代价），用分布式更优响应迭代，验证收敛到纯 Nash；讨论它与匈牙利算法（集中式最优分配）的差异。
  \item 论证：(a) 为什么"所有 agent 共享同一团队奖励"让博弈成为势博弈（团队奖励 $=$ 势函数）；(b) 这如何解释 QMIX/VDN 等合作 MARL 的收敛性；(c) 当 agent 奖励\emph{不完全一致}（混合合作-竞争）时，势博弈保证为何失效，该怎么办？
\end{enumerate}
\end{practice}

% ════════════════════════════════════════════════════════════════
\section{多智能体强化学习的形式化：Markov Game 与 Dec-POMDP}
\label{sec:sm-marl-form}

前两节用"显式建模 $+$ 数值求解"处理博弈（CBF-QP、势函数梯度）。但当 agent 很多、博弈很大（扑克、围棋、大规模 swarm），或代价函数本身未知/难以手工设计时，另一条路是\textbf{用强化学习让 agent 在反复对弈（自博弈）中学出均衡}。要走这条路，先得有它的公共语言——正如单智能体强化学习有马尔可夫决策过程（MDP）作为公共语言。

\paragraph{动机：单智能体 MDP 为什么不够用。}
单智能体强化学习的世界观里，环境用 MDP $(\mathcal{S},\mathcal{A},P,R,\gamma)$ 描述，核心假设是\textbf{环境是被动的、固定的}——转移核 $P$ 和奖励 $R$ 不会因为智能体学到了什么而改变。正是这个假设，让贝尔曼算子成为压缩映射（\cref{thm:plan-bellman} 的收敛根基）、让 Q-learning 收敛、让经验回放里的样本可当独立同分布用。现在考虑一个真实多机器人场景：仓库里 4 台 AGV 协同运货、互不相撞。如果硬套单智能体 MDP，会立刻遇到尴尬：\textbf{对 1 号 AGV 来说，"环境"是什么？} 它的下一状态不仅取决于自己的动作，还取决于另外 3 台的动作——而那 3 台也在各自学习、各自改变策略。于是 1 号眼中的"环境"会变：今天 2 号喜欢走左边，1 号学会走右边；明天 2 号改走右边，1 号昨天的策略就失效了。单智能体 MDP \textbf{无法表达"环境里有其他正在学习的决策者"}。我们需要把"多个决策者同时在共享环境里行动、每个人的动作影响所有人"显式写进数学结构。

\paragraph{反面：硬套单智能体的两种做法都踩坑。}
\begin{itemize}
  \item \textbf{独立学习（Independent Learning）}：给每个智能体配一份独立的单智能体算法（每人一份 DQN 或 PPO，即 IQL/IPPO），把其他智能体都当"环境的一部分"。致命问题：从 $i$ 的视角，有效转移核依赖他人策略 $\{\pi_j\}_{j\neq i}$，他人策略在训练中变化 $\Rightarrow$ 有效核漂移 $\Rightarrow$ 环境不再是固定 MDP（非平稳，\cref{sec:sm-nonstationary} 详述）。
  \item \textbf{联合学习（Joint Action Learning）}：把所有智能体捏成一个超级智能体，动作是联合动作 $\mathbf{a}=(a_1,\dots,a_N)$，策略 $\pi(\mathbf{a}\mid s)$ 直接在联合动作空间决策。这回到了良定义的平稳 MDP，但代价是\emph{联合动作空间随智能体数指数爆炸}（$|\mathcal{A}|^N$，10 个智能体各 5 动作就是 $\sim10^7$），且部署时超级智能体需同时拿到所有观测——\emph{执行时不可分布式实现}。
\end{itemize}

\begin{insight}
独立学习与联合学习是两个极端，各踩一坑。独立学习保住"分布式执行"（每人只看自己）却丢了"平稳性"（环境在漂）；联合学习保住"平稳性"（回到良定义 MDP）却丢了"可扩展性"和"分布式执行"。MARL 的全部算法设计，本质上都是在这两个极端之间寻找一个\textbf{既保平稳性、又能分布式执行、还不指数爆炸}的中间地带——这个中间地带叫 CTDE（\cref{sec:sm-ctde}）。但要精确描述它，得先有能同时刻画"联合"与"分布式"的形式化语言。
\end{insight}

\paragraph{历史：从博弈论到 Dec-POMDP。}
MARL 的形式化有两个源头。\emph{博弈论源头}：Shapley 1953 在随机环境下提出随机博弈（Stochastic Game）；Littman 把它引入强化学习、称 Markov Game，并给出零和两人博弈的 minimax-Q\cite{littman1994markov}。这一支强调\textbf{博弈结构}——智能体可能合作、竞争、混合，要找的是某种\emph{均衡}而非单一最优。\emph{运筹学/控制源头}：Bernstein 等形式化了去中心化部分可观测 MDP（Dec-POMDP）\cite{bernstein2002decpomdp}——所有智能体共享同一奖励（纯合作），但每个只观测全局状态的一部分、执行时不能通信；他们证明了一个令人沮丧的结果：\textbf{有限时域 Dec-POMDP 的最优求解是 NEXP-完全的}（比单智能体 POMDP 的 PSPACE 还难一个数量级），这解释了为什么所有实用方法都转向"用集中信息辅助训练、再蒸馏成分布式策略"（即 CTDE）。两个源头在合作 MARL 里汇合：\textbf{合作型随机博弈 $=$ 状态完全可观测的 Dec-POMDP}。

\subsection{随机博弈（Markov Game）}

\begin{definition}[随机博弈 / Markov Game]\label{def:sm-markovgame}
随机博弈是一个元组 $\mathcal{G}=\langle N,\ \mathcal{S},\ \{\mathcal{A}_i\}_{i=1}^N,\ P,\ \{r_i\}_{i=1}^N,\ \gamma\rangle$，其中 $N$ 是智能体数，$\mathcal{S}$ 是全局状态空间，$\mathcal{A}_i$ 是智能体 $i$ 的动作空间（联合动作空间 $\boldsymbol{\mathcal{A}}=\mathcal{A}_1\times\dots\times\mathcal{A}_N$，联合动作 $\mathbf{a}=(a_1,\dots,a_N)$），$P:\mathcal{S}\times\boldsymbol{\mathcal{A}}\to\Delta(\mathcal{S})$ 是\textbf{联合转移核}（下一状态取决于当前状态和\emph{所有人的}联合动作），$r_i:\mathcal{S}\times\boldsymbol{\mathcal{A}}\to\mathbb{R}$ 是智能体 $i$ 的奖励函数（\emph{别人的动作会影响我的奖励}），$\gamma\in[0,1)$ 是折扣因子。
\end{definition}

每个智能体 $i$ 有策略 $\pi_i(a_i\mid s)$（先设完全可观测），联合策略 $\boldsymbol{\pi}=(\pi_1,\dots,\pi_N)$。智能体 $i$ 的值函数
\begin{equation}
  V_i^{\boldsymbol{\pi}}(s)=\mathbb{E}\Big[\sum_{t=0}^\infty\gamma^t\,r_i(s_t,\mathbf{a}_t)\ \Big|\ s_0=s,\ \mathbf{a}_t\sim\boldsymbol{\pi}(\cdot\mid s_t)\Big].
  \label{eq:sm-mg-value}
\end{equation}
\textbf{关键观察}：$V_i$ 带的是上标 $\boldsymbol{\pi}$ 而不只是 $\pi_i$——智能体 $i$ 的回报不仅取决于自己的策略，还取决于所有其他智能体的策略 $\boldsymbol{\pi}_{-i}$。所以"$i$ 的最优策略"在博弈里没有绝对意义，只能说"给定别人的策略，$i$ 的最优响应（best response）" $\pi_i^{\text{BR}}(\boldsymbol{\pi}_{-i})=\arg\max_{\pi_i}V_i^{(\pi_i,\boldsymbol{\pi}_{-i})}(s)$。当所有智能体的策略互为最优响应——没人能单方面改变策略来提高回报——这组策略构成一个\textbf{纳什均衡}：
\begin{equation}
  \forall i,\quad V_i^{\boldsymbol{\pi}^*}(s)\ \ge\ V_i^{(\pi_i,\,\boldsymbol{\pi}^*_{-i})}(s)\quad\text{对所有 }\pi_i,\,s.
  \label{eq:sm-nash}
\end{equation}

\begin{note}
\textbf{纳什均衡 $\neq$ 集体最优.}\;
随机博弈之于 MDP，如同纳什均衡之于单变量最优化：单变量最优化找一个点让一个函数最大，纳什均衡找一组策略让 $N$ 个函数\emph{同时各自对自己的变量最优}。但关键不同：单变量最优化的全局最优唯一且明确"更好"；纳什均衡可能有多个、不同均衡不一定可比（一个对智能体 1 好、对 2 差），甚至均衡的总收益可能远不如某个非均衡的协作点（囚徒困境：均衡是"双方背叛"，但"双方合作"才集体最优、却不是均衡）。\textbf{不要把"找到纳什均衡"等同于"找到最优"}——这是 MARL 初学者最常见的误解。
\end{note}

按奖励结构，随机博弈分三类（MARL 的顶层分类）：\textbf{合作（cooperative）}——所有智能体共享同一奖励 $r_1=\dots=r_N=r$（多机协同搬运、编队、覆盖）；\textbf{竞争（competitive）}——零和 $\sum_i r_i=0$（对抗、围捕-逃逸）；\textbf{混合（mixed）}——一般和 general-sum、无约束（自动驾驶交汇、市场、社会困境）。本章主战场是\textbf{合作}——绝大多数多机器人协作任务的结构。合作场景有一个巨大简化：因为所有人共享奖励 $r$，所有人的值函数相同 $V_1^{\boldsymbol{\pi}}=\dots=V_N^{\boldsymbol{\pi}}=V^{\boldsymbol{\pi}}$，于是"找纳什均衡"退化为"找联合策略最大化共享回报"——一个\textbf{良定义的单目标优化}，没有了"我好你坏"的均衡选择难题。换句话说，\textbf{合作 MARL 在目标层面回到了单目标优化，困难只剩"信息分散"和"非平稳性"}——这正是 CTDE 能发挥作用的根本原因。

\subsection{Dec-POMDP（合作 \texorpdfstring{$+$}{+} 部分可观测）}

随机博弈假设每个智能体能看到全局状态 $s$。但真实多机器人系统里，每个机器人只能看到自己的传感器——局部、带噪声、视野受限。把"部分可观测"加进合作随机博弈，就得到本章最核心的形式化。

\begin{definition}[去中心化部分可观测 MDP（Dec-POMDP）]\label{def:sm-decpomdp}
Dec-POMDP 是一个元组 $\mathcal{M}=\langle N,\ \mathcal{S},\ \{\mathcal{A}_i\},\ P,\ r,\ \{\mathcal{O}_i\},\ \{O_i\},\ \gamma\rangle$，在随机博弈基础上：\textbf{共享奖励} $r:\mathcal{S}\times\boldsymbol{\mathcal{A}}\to\mathbb{R}$（纯合作，只有一个团队奖励，下标 $i$ 没了）；$\mathcal{O}_i$ 是智能体 $i$ 的观测空间；$O_i:\mathcal{S}\times\boldsymbol{\mathcal{A}}\to\Delta(\mathcal{O}_i)$ 是观测函数；\textbf{关键约束：每个智能体的策略只能依赖自己的局部历史}。
\end{definition}

最后一条是 Dec-POMDP 的灵魂。因为部分可观测，单步观测 $o_i$ 一般不是马尔可夫的（不同历史可能给出相同 $o_i$ 却需要不同动作——状态混叠 state aliasing）。所以智能体 $i$ 的策略不能只看当前观测，而要看\textbf{局部动作-观测历史} $\tau_i=(o_i^0,a_i^0,o_i^1,a_i^1,\dots,o_i^t)$，策略写成 $\pi_i(a_i\mid\tau_i)$。注意 $\tau_i$ 是智能体 $i$\emph{自己}的历史，不含别人的观测或动作——这就是"去中心化"的含义。目标是找联合策略最大化共享期望回报：
\begin{equation}
  \boldsymbol{\pi}^*=\arg\max_{\boldsymbol{\pi}}\ \mathbb{E}\Big[\sum_{t=0}^\infty\gamma^t\,r(s_t,\mathbf{a}_t)\ \Big|\ a_i^t\sim\pi_i(\cdot\mid\tau_i^t)\Big].
  \label{eq:sm-decpomdp-obj}
\end{equation}
拆开看每部分的工程含义：目标里是共享奖励 $r$、求和不带 $i$——这是"合作"，所有智能体为同一团队目标服务（也是后面信用分配问题的来源，\cref{sec:sm-credit}）；约束 $a_i^t\sim\pi_i(\cdot\mid\tau_i^t)$——这是"去中心化执行"，每个 $\pi_i$ 只吃自己的历史（部署时不需中央协调器、不需通信）；$\arg\max$ 对联合策略 $\boldsymbol{\pi}$——优化是联合的，但解出来是 $N$ 个独立的 $\pi_i$。

\begin{insight}
Dec-POMDP 把 MARL 的难度精确定位成一句话——\textbf{"用只能看局部的策略，去最大化一个全局的团队目标"}。如果允许每个 $\pi_i$ 看全局状态 $s$，问题退化成单智能体 MDP；如果允许执行时自由通信、共享所有观测，问题退化成集中式。Dec-POMDP 的全部难度，来自"优化是全局的、但执行必须用局部信息"这个不可调和的张力。NEXP-完全的复杂度，就是这个张力的代价：每个智能体不仅要推断真实状态，还要推断别人会怎么做——而别人也在推断我会怎么做，形成无穷递归；且执行时不能通信，意味着智能体必须在训练阶段就把"如何协调"编码进各自策略，让它们在执行时即使不交流也能默契配合。
\end{insight}

用一张表把四个形式化的关系钉死（读法：从 MDP 出发"加料"）：

\begin{table}[H]\centering\small
\caption{四套形式化的关系}
\label{tab:sm-frameworks}
\begin{tabular}{lcccl}
\toprule
框架 & 智能体数 & 可观测性 & 奖励结构 & 退化关系 \\
\midrule
MDP & $1$ & 完全 & 单一 & 基准 \\
POMDP & $1$ & 部分 & 单一 & MDP $+$ 部分可观测 \\
Markov Game & $N$ & 完全 & 一般和 & MDP $+$ 多智能体（一般奖励） \\
Dec-POMDP & $N$ & 部分 & 共享（合作） & Markov Game $+$ 部分可观测 $+$ 强制合作 \\
\bottomrule
\end{tabular}
\end{table}

\begin{pitfall}[概念误区：MARL 就是把单智能体 RL 复制 $N$ 份]
很多人以为"MARL 就是把单智能体 RL 复制 $N$ 份"。\textbf{不是}。若真只是复制 $N$ 份独立运行，那叫独立学习（IQL/IPPO），它会被非平稳性击垮（\cref{sec:sm-nonstationary}）。MARL 的真正困难在于这 $N$ 份必须\emph{协调}——它们共享一个状态空间、动作互相影响、奖励互相耦合。\textbf{MARL 不是"$N$ 个单智能体 RL"，而是"一个定义在联合空间上、但解必须能分解成 $N$ 个局部策略的优化问题"}。另一个常见陷阱：把全局状态 $s$ 等同于观测拼接 $(o_1,\dots,o_N)$——两者通常\emph{不等}（每个 $o_i$ 都是部分的、可能重叠、可能有盲区），只在"所有观测无重叠且并起来覆盖整个状态"的特殊情况才相等。实现集中式 critic 时这个区分很关键（MAPPO 的 \texttt{cent\_obs} 与 actor 的 \texttt{obs} 严格分开）。
\end{pitfall}

\begin{practice}
\begin{enumerate}
  \item 把"$N$ 台 AGV 在仓库取送货、不能相撞"写成 Dec-POMDP，逐项给出 $\mathcal{S},\{\mathcal{A}_i\},P,r,\{\mathcal{O}_i\},\{O_i\}$；回答：(a) 若允许 AGV 实时通信共享位置，会退化成什么框架？(b) 多机路径规划（MAPF）可看成它的什么特例（提示：奖励/动力学设成什么就变回离散无碰撞路径搜索）？
  \item 证明：共享奖励的合作随机博弈中所有智能体的值函数相等 $V_1^{\boldsymbol{\pi}}=\dots=V_N^{\boldsymbol{\pi}}$；论证为什么这使"找纳什均衡"退化为单目标优化；给出一个两智能体反例说明非共享奖励（一般和）下此简化失效。
  \item 设计一个最小 $2\times2$ 网格、2 智能体合作任务，使得用单帧观测 $\pi_i(a_i\mid o_i)$ 的最优策略严格劣于用历史 $\pi_i(a_i\mid\tau_i)$ 的；明确指出哪两个真实状态被单帧观测混叠了。
\end{enumerate}
\end{practice}

% ════════════════════════════════════════════════════════════════
\section{非平稳性：MARL 的第一性困难}
\label{sec:sm-nonstationary}

\paragraph{动机：为什么"复制 $N$ 份单智能体 RL"必然失败。}
最自然的 MARL 想法是独立学习：给每个智能体配一份成熟算法、各自训练。这个想法如此自然，几乎每个人入门 MARL 时都会先试一遍——然后撞墙。实际现象是：训练曲线剧烈震荡、不收敛，或收敛到很差的策略；换学习率、换网络、加大样本量都没用。问题的根源不是超参、不是实现 bug，而是\textbf{框架性的}——单智能体 RL 的全部理论保证建立在一个被多智能体悄悄违反的假设上。这一节把它精确写出来。

\subsection{非平稳性的精确定义与三重破坏}

\paragraph{第一步：写出单个智能体看到的"有效环境"。}
在 Dec-POMDP 里，真实的联合转移核 $P(s'\mid s,\mathbf{a})$ 是固定的——这一点没变。但从智能体 $i$ 的视角，它看不到也控制不了别人的动作 $\mathbf{a}_{-i}$，别人的动作按各自策略 $\pi_j$ 采样。所以 $i$ 实际面对的"有效转移核"是把别人的动作积分掉之后的边缘化核：
\begin{equation}
  P_i^{\boldsymbol{\pi}_{-i}}(s'\mid s,a_i)=\sum_{\mathbf{a}_{-i}}P(s'\mid s,a_i,\mathbf{a}_{-i})\prod_{j\neq i}\pi_j(a_j\mid s),
  \label{eq:sm-effective-core}
\end{equation}
有效奖励同理被别人的策略平均掉：$R_i^{\boldsymbol{\pi}_{-i}}(s,a_i)=\sum_{\mathbf{a}_{-i}}r(s,a_i,\mathbf{a}_{-i})\prod_{j\neq i}\pi_j(a_j\mid s)$。

\paragraph{第二步：看出问题。}
这两个有效核里都带着 $\boldsymbol{\pi}_{-i}$——别人的策略。训练中别人也在学、在变，$\boldsymbol{\pi}_{-i}$ 随训练步 $k$ 变化，记作 $\boldsymbol{\pi}_{-i}^{(k)}$。于是 $i$ 面对的不是一个固定 MDP，而是一串随时间变化的 MDP：$(P_i^{(0)},R_i^{(0)})\to(P_i^{(1)},R_i^{(1)})\to\cdots$。

\begin{definition}[非平稳性 non-stationarity]\label{def:sm-nonstationarity}
从单个智能体的视角，由于其他智能体的策略在学习过程中不断变化，该智能体所面对的有效转移核 \cref{eq:sm-effective-core} 与有效奖励随时间漂移，使其环境不再是一个固定的 MDP，而是一个非平稳的决策过程。
\end{definition}

\paragraph{破坏一：压缩映射失效（理论收敛保证崩塌）。}
回顾单智能体 Q-learning 为什么收敛（\cref{thm:plan-bellman} 的同一逻辑）：贝尔曼最优算子 $(\mathcal{T}Q)(s,a)=R(s,a)+\gamma\sum_{s'}P(s'\mid s,a)\max_{a'}Q(s',a')$ 是 $\gamma$-压缩映射——$\|\mathcal{T}Q_1-\mathcal{T}Q_2\|_\infty\le\gamma\|Q_1-Q_2\|_\infty$。由 Banach 不动点定理，迭代 $Q\leftarrow\mathcal{T}Q$ 收敛到唯一不动点 $Q^*$。这个证明的\emph{每一步都用到 $P,R$ 固定}。现在把 $P,R$ 换成随时间变的 $P_i^{(k)},R_i^{(k)}$，智能体 $i$ 的贝尔曼算子变成 $\mathcal{T}^{(k)}$——一个\emph{会变的算子}。即使每个 $\mathcal{T}^{(k)}$ 单独看是压缩的，它们的不动点 $Q_i^{*(k)}$ 也是个移动靶：$i$ 朝 $Q_i^{*(0)}$ 努力到一半，目标已漂移到 $Q_i^{*(1)}$ 了。\textbf{追逐一个移动靶，迭代不保证收敛}——这就是独立 Q-learning（IQL）没有收敛保证的数学根源，不是它实现得不好，而是它要解的根本不是一个不动点问题。

\begin{insight}
单智能体 RL 之所以能收敛，是因为它在和一个\textbf{死的}环境博弈——环境固定，智能体单方面优化，朝固定靶逼近。MARL 之所以难，是因为每个智能体都在和一群\textbf{活的、也在学习的}对手/队友博弈——靶随着所有人的学习一起动。\textbf{非平稳性的本质，就是"固定靶"退化为"移动靶"}。后面 CTDE 的全部精妙，在于它想办法把移动靶重新"钉住"——给 critic 喂别人的动作，相当于告诉它"靶现在在哪"，从而局部恢复平稳性。
\end{insight}

\paragraph{破坏二：经验回放失效（off-policy 数据过期）。}
DQN 的样本效率靠经验回放——把历史转移存进 buffer 反复采样，打破时间相关性、近似独立同分布。单智能体下成立，因为环境固定：三年前的转移和今天的描述同一个 $P,R$。多智能体下，buffer 里一条旧转移是在"别人用旧策略 $\boldsymbol{\pi}_{-i}^{(k_{\text{old}})}$"时产生的，但现在别人策略已变成 $\boldsymbol{\pi}_{-i}^{(k_{\text{now}})}$，这条旧转移描述的有效核 $P_i^{(k_{\text{old}})}$ 和当前的 $P_i^{(k_{\text{now}})}$ 不是同一分布。\textbf{用过期的、来自不同环境的数据训练当前策略，会引入系统性偏差}——越大的 buffer、越久的历史，过期越严重。这解释了为什么 on-policy 的 MAPPO（每次用最新数据、用完即弃）在 MARL 里出奇地稳——它根本不存 buffer，自然没有数据过期。这也是 MADDPG 必须把别人的动作 $\mathbf{a}_{-i}$ 存进 buffer、作为 critic 输入的动机（\cref{sec:sm-policygrad}）：条件给定了别人的动作，值函数就不依赖别人的策略了，旧转移依然有效。

\paragraph{破坏三：探索-利用的耦合。}
你的探索性动作会改变别人看到的环境，触发别人改变策略，反过来又改变你的环境。两个智能体可能陷入"你变我变、互相追逐"的振荡（博弈论里"最优响应动力学不收敛"）。MARL 的探索不是 $N$ 个独立探索的简单叠加，而是耦合的。

\subsection{系统性分类：MARL 的四大困难}
\label{sec:sm-four-difficulties}

非平稳性是第一性困难，但不是唯一。把 MARL 的困难做穷举式维度分类，你就有了一个诊断框架：拿到任何方法，都能问"它在攻击四大困难中的哪个、用什么手段"。

\begin{table}[H]\centering\small
\caption{MARL 的四大核心困难}
\label{tab:sm-four-diff}
\begin{tabular}{p{0.16\textwidth} p{0.24\textwidth} p{0.24\textwidth} p{0.24\textwidth}}
\toprule
困难 & 一句话定义 & 主要应对手段 & 本章对应 \\
\midrule
\textbf{① 非平稳性} & 别人在学，我的环境在漂 & 集中式 critic（喂别人动作/全局状态）$\to$ 局部恢复平稳 & \cref{sec:sm-ctde,sec:sm-policygrad} \\
\textbf{② 信用分配} & 团队拿了奖励，谁的功劳？ & 差分奖励、反事实基线、值分解 & \cref{sec:sm-credit,sec:sm-valuedecomp} \\
\textbf{③ 维度爆炸} & 联合动作空间随 $N$ 指数增长 & 分布式执行（各自 argmax）、参数共享、分解 & \cref{sec:sm-ctde,sec:sm-valuedecomp} \\
\textbf{④ 均衡与协调} & 多个均衡选哪个？不通信如何默契？ & 合作时退化为单目标；协调靠训练期对齐 & \cref{sec:sm-marl-form,sec:sm-policygrad} \\
\bottomrule
\end{tabular}
\end{table}

\textbf{② 信用分配}是合作 MARL 特有的、和非平稳性并列的第二大困难。合作场景只有一个团队奖励 $r$。当团队成功拿到 $+100$，策略梯度要更新每个智能体——但这 $+100$ 里有多少是 1 号的贡献、多少是 2 号的？独立学习者 $i$ 用团队奖励做策略梯度 $\nabla_{\theta_i}J\approx\mathbb{E}[\nabla_{\theta_i}\log\pi_i(a_i\mid o_i)\cdot r(s,\mathbf{a})]$，问题是 $r(s,\mathbf{a})$ 里含 $\mathbf{a}_{-i}$ 的贡献——这部分对 $a_i$ 是纯噪声、方差很大，会淹没 $a_i$ 自己的微弱信号。信用分配的本质，就是从团队奖励里把"我的那一份"剥离出来（\cref{sec:sm-credit}）。

\textbf{③ 维度爆炸}：联合动作空间 $|\boldsymbol{\mathcal{A}}|=\prod_i|\mathcal{A}_i|$ 随 $N$ 指数增长，有学习层面（联合 $Q$ 的输出空间爆炸）和决策层面（$\arg\max_{\mathbf{a}}Q$ 搜不动且需全局信息）两个危害。应对核心是\emph{分解 $+$ 分布式}——学 $N$ 个小的个体 utility $Q_i(o_i,a_i)$，让每个智能体独立 $\arg\max_{a_i}Q_i$（线性而非指数）。但这里有陷阱：\textbf{分布式的局部 argmax 凑起来凭什么等于集中式的全局 argmax}？这正是 \cref{sec:sm-valuedecomp} 的 IGM 原则要解决的。

\textbf{④ 均衡选择与协调}：一般和博弈可能有多个纳什均衡，智能体如何收敛到同一个、且"好"的那个。合作场景在这一点上有巨大简化（共享奖励 $\Rightarrow$ 单目标），但\emph{协调}问题依然存在——即使目标一致，$N$ 个智能体执行时不能通信，如何保证局部决策拼成协调的整体？HAPPO 用"顺序更新"缓解（\cref{sec:sm-policygrad}）。

\begin{pitfall}[概念误区：非平稳性可以靠"调小学习率"解决]
把独立学习不收敛归因于学习率太大、试图减小各智能体学习率来稳定训练，是常见误区。减小学习率确实让震荡看起来平缓，但收敛速度急剧下降，且最终仍收敛不到好策略——本质问题没解决，只是把"快速震荡"变成"缓慢漂移"。根本原因：非平稳性是\emph{框架性}问题（追逐移动靶），不是优化器步长问题。调小学习率只是让所有智能体都走得慢，移动靶和追逐者一起慢下来，相对运动没变。真正的解药是\emph{改变信息结构}——CTDE 让 critic 看到别人的动作/全局状态，把移动靶钉住（\cref{sec:sm-ctde}），而非改变优化超参。另一个相关陷阱：在多智能体 off-policy 训练中照搬单智能体的大 replay buffer——buffer 越大、过期数据占比越高、偏差越严重；off-policy 路线（MADDPG）必须把联合动作存进 buffer 并作为集中 critic 输入，或改用 on-policy（MAPPO）从源头避开。
\end{pitfall}

\begin{practice}
\begin{enumerate}
  \item 取一个具体的 $2\times2$ 团队奖励矩阵，让两个智能体各自跑独立 Q-learning，手动模拟 3--4 轮"最优响应"（固定 B 求 A 的有效奖励与最优动作，再固定 A 求 B……），观察是否出现振荡，用这个手算例子具体说明"移动靶"是怎么产生的。
  \item 设智能体 $i$ 的 replay buffer 容量 $M$，每步采集 1 条新转移、淘汰 1 条最旧；假设他人策略每 $T$ 步发生一次"显著变化"。论证 buffer 里"来自当前有效环境"的转移比例约为 $\min(1,T/M)$，据此说明为什么他人策略快速变化（$T$ 小）时应当用更小的 buffer，以及为什么 on-policy 方法天然规避此问题。
  \item 跳读 \cref{sec:sm-valuedecomp} 的 QMIX 或 \cref{sec:sm-policygrad} 的 MAPPO，不看 \cref{tab:sm-four-diff}，自己判断它主要攻击四大困难中的哪个、用了什么机制，再回来对照。
\end{enumerate}
\end{practice}

% ════════════════════════════════════════════════════════════════
\section{CTDE：集中训练，分布执行}
\label{sec:sm-ctde}

\paragraph{动机：在两个极端之间找中间道路。}
把 \cref{sec:sm-nonstationary} 的结论摆在一起，矛盾就清楚了。我们想要一个 MARL 方案，同时满足三个看似冲突的要求：(1) \textbf{执行时分布式}——部署到真实多机器人时，每个机器人只能用自己的局部观测 $o_i$ 决策，不能依赖全局状态、不能假设实时通信（Dec-POMDP 的硬约束）；(2) \textbf{训练时不被非平稳性击垮}；(3) \textbf{可扩展}——不能随智能体数指数爆炸。独立学习满足 1、3 违反 2，联合学习满足 2 违反 1、3。关键洞察来自一个\textbf{不对称}——\emph{训练和执行是两个不同阶段，约束不一样}：执行阶段在真实系统里约束严格（只有局部信息、不能通信）；训练阶段通常在\emph{仿真器}里，约束宽松（仿真器是上帝视角，知道所有智能体的精确状态、所有人的动作、环境全部真值，这些信息\emph{唾手可得且免费}）。既然如此，为什么不\textbf{在训练时大方地用全局信息、在执行时只保留局部信息}？这就是 CTDE 的核心思想。

\begin{definition}[集中训练分布执行 CTDE]\label{def:sm-ctde}
CTDE（Centralized Training, Decentralized Execution）是一类 MARL 算法范式：\textbf{训练时}使用一个或多个\emph{集中式组件}（critic、混合网络、联合值函数等），可访问全局信息——全局状态 $s$、所有智能体的观测 $\{o_i\}$、所有智能体的动作 $\{a_i\}$；\textbf{执行时}只保留\emph{分布式组件}（每个智能体的 actor / 个体值函数 $\pi_i(a_i\mid o_i)$ 或 $Q_i(o_i,a_i)$），只依赖本地观测 $o_i$（或本地历史 $\tau_i$），不需要全局信息、不需要智能体间通信。
\end{definition}

\paragraph{理论：为什么集中式 critic 能治非平稳性。}
回顾 \cref{sec:sm-nonstationary}：独立学习者 $i$ 的值估计依赖他人策略 $\boldsymbol{\pi}_{-i}$（有效核 \cref{eq:sm-effective-core} 随之漂移）。CTDE 的集中式 critic 不再估"边缘化掉别人之后"的值，而是估\textbf{以全局状态/联合动作为条件}的值 $Q(s,a_1,\dots,a_N)$ 或 $V(s)$。关键在于：一旦 critic 的输入里\emph{显式包含了别人的动作 $\mathbf{a}_{-i}$（或全局状态 $s$，它编码了别人的位形）}，别人策略变化的影响就被"条件"吸收了。对一个固定的 $(s,\mathbf{a})$，转移 $P(s'\mid s,\mathbf{a})$ 是\emph{真正固定的}（这是 Dec-POMDP 里没变的那一条）。所以以联合动作为条件的 $Q(s,\mathbf{a})$ 满足一个良定义、固定的贝尔曼方程：
\begin{equation}
  Q(s,\mathbf{a})=r(s,\mathbf{a})+\gamma\,\mathbb{E}_{s'\sim P(\cdot\mid s,\mathbf{a})}\Big[\mathbb{E}_{\mathbf{a}'\sim\boldsymbol{\pi}(\cdot\mid s')}\big[Q(s',\mathbf{a}')\big]\Big].
  \label{eq:sm-central-bellman}
\end{equation}
这个 $Q(s,\mathbf{a})$ 的贝尔曼算子相对于固定核 $P$ 是压缩的——\textbf{平稳性在"联合动作为条件"这个层面被恢复了}。别人策略的变化只通过下一步动作分布 $\boldsymbol{\pi}(\cdot\mid s')$ 进入（on-policy 评估里本来就有的、可控的部分），而不再让转移核本身漂移。

\begin{insight}
CTDE 治非平稳性的机理可以用一句话概括——\textbf{"把漂移的东西从'环境'里挪到'条件'里"}。独立学习里，别人的动作藏在被边缘化掉的有效核里，看不见摸不着、随别人策略漂移；CTDE 里，别人的动作被显式拎出来当 critic 的输入条件，于是"以联合动作为条件"的那个值函数是相对固定核定义的、平稳的。漂移没有消失，但它从"破坏压缩性的环境漂移"降级为"评估目标里可控的策略项"。这是整个 MARL 能工作的数学支点。
\end{insight}

\textbf{为什么分布式 actor 不破坏这个好处？} 因为 critic 只在\emph{训练时}用来给 actor 提供梯度信号（优势/基线）。actor 自己始终是分布式的 $\pi_i(a_i\mid o_i)$，它从集中式 critic 那里"借"全局视角来改进，但自己只学会用局部观测做决策。训练结束，critic 这个"脚手架"拆掉，留下的 actor 天然满足分布式执行。这就像考前可以查所有资料复习（集中训练），但考场上只能靠脑子（分布执行）——复习用的资料把知识内化进了脑子，考场上不需要再带资料。

CTDE 是一个范式、不是单一算法。它的"集中式组件"可以有不同形态，对应不同算法家族，把这三种变体梳理清楚就有了 MARL 算法的"分类骨架"：

\begin{table}[H]\centering\small
\caption{CTDE 的三种工程变体}
\label{tab:sm-ctde-variants}
\begin{tabular}{p{0.16\textwidth} p{0.34\textwidth} p{0.22\textwidth} p{0.16\textwidth}}
\toprule
变体 & 集中式组件是什么 & 代表算法 & 适用 \\
\midrule
\textbf{集中式 critic} & 看全局的 critic $Q(s,\mathbf{a})$ 或 $V(s)$，给分布式 actor 提供优势/基线 & MADDPG, COMA, MAPPO, HAPPO & 策略梯度路线；连续/离散均可 \\
\textbf{值分解} & 混合网络 $f$ 把个体 utility 聚合成 $Q_{\text{tot}}=f(Q_1,\dots,Q_N;s)$，只在训练时存在 & VDN, QMIX, QTRAN, QPLEX & 值函数路线；主要离散动作 \\
\textbf{可微通信} & 训练时端到端反传的通信通道或共享表征 & RIAL/DIAL, CommNet & 需学习通信协议的任务 \\
\bottomrule
\end{tabular}
\end{table}

前两种是本章主线（\cref{sec:sm-valuedecomp} 值分解、\cref{sec:sm-policygrad} 集中 critic）。

\begin{note}
\textbf{CTDE 的边界与代价.}\;它不是银弹。(1) \emph{不能产生执行时的实时协调}——CTDE 的协调是"离线编译"进各自策略的，执行时各自独立跑、不再沟通；若任务要求执行时临场协商，需叠加显式通信。(2) \emph{集中式 critic 的输入随 $N$ 增长}（线性，不是指数，但也增长），$N$ 很大时难训（Mean-Field MARL 进一步处理）。(3) \emph{"全局状态 $s$"的获取假设}——纯仿真里没问题，真实数据离线训练时只能用观测拼接近似（\cref{sec:sm-marl-form} 已警示两者不等）。(4) \emph{训练-执行的分布偏移}——执行时 actor 独立决策产生的联合动作分布，可能和训练时 critic 见过的不同。
\end{note}

\begin{note}
\textbf{CTDE（信息结构）$\neq$ 参数共享（权重复用）.}\;参数共享（\cref{sec:sm-policygrad} 的 MAPPO 会用）说的是"$N$ 个 actor 用同一套权重"（省参数、提样本效率）；CTDE 说的是"训练时有个看全局的集中组件"（治非平稳性）。两者\emph{正交}：一个算法可以只 CTDE 不共享参数（如 MADDPG 每个 agent 独立 actor），也可以既 CTDE 又共享参数（如 MAPPO）。读论文时务必把这两个概念分开，否则会把 MAPPO 的两个独立设计（全局 critic $+$ 参数共享）混成一个。
\end{note}

\begin{pitfall}[思维陷阱：有了 CTDE 就自动解决了信用分配]
认为"用了集中式 critic / CTDE，信用分配问题就自动消失"是常见误区——用了 MAPPO（CTDE）却发现智能体仍有搭便车行为。根本原因：CTDE 主要解决的是\emph{非平稳性}和\emph{维度爆炸（决策层面）}，信用分配是一个\emph{正交的}困难（\cref{tab:sm-four-diff}）。一个共享团队奖励、用普通 $V(s)$ 当 critic 的 CTDE 方法，其优势 $A=r+\gamma V(s')-V(s)$ 对所有智能体一样，并没有区分各自贡献——信用分配靠额外机制（反事实基线 COMA、值分解 QMIX）。把"治非平稳性"和"解信用分配"当两件事分别处理。还有一个相关误区：以为训练好的集中 critic 在部署时也要运行，于是担心"真实系统拿不到全局状态怎么办"——集中 critic 是"教练"不是"运动员"，训练后整个丢弃，部署时只跑分布式 actor，完全满足分布式执行。
\end{pitfall}

\begin{practice}
\begin{enumerate}
  \item 对"3 无人机覆盖"任务设计：(a) 集中式 critic 的输入应含哪些信息（全局状态哪些成分 $+$ 是否需联合动作）？(b) 分布式 actor 的输入是什么？(c) 若仿真器只暴露各无人机观测、不暴露真正全局状态，如何近似构造 critic 输入？这个近似会损失什么？
  \item 写出以联合动作为条件的 $Q(s,\mathbf{a})$ 的贝尔曼方程 \cref{eq:sm-central-bellman}，论证相对固定核 $P(s'\mid s,\mathbf{a})$ 它的算子是 $\gamma$-压缩的；说明为什么独立学习者的 $Q_i(o_i,a_i)$ 对应的算子\emph{不是}相对固定核的压缩。
  \item 给出四个组合各一个具体算法或设计：(a) CTDE $+$ 参数共享；(b) CTDE $+$ 不共享参数；(c) 非 CTDE $+$ 参数共享；(d) 非 CTDE $+$ 不共享参数；对每个说明它在"治非平稳性"和"样本效率"两维度的位置。
\end{enumerate}
\end{practice}

% ════════════════════════════════════════════════════════════════
\section{值分解：VDN、QMIX 与 IGM 原则}
\label{sec:sm-valuedecomp}

CTDE 给了骨架，现在填第一种"集中式组件"——值分解（value decomposition）。核心矛盾：我们想让每个智能体执行时独立地 $\arg\max_{a_i}Q_i(o_i,a_i)$（治维度爆炸），但又想让这些局部贪婪凑起来恰好等于集中式的全局最优 $\arg\max_{\mathbf{a}}Q_{\text{tot}}(s,\mathbf{a})$。

\paragraph{动机：分布式 argmax 凭什么等于集中式 argmax。}
回到维度爆炸（\cref{sec:sm-four-difficulties} 困难③）。最自然的想法是给每个智能体 $i$ 学一个\textbf{个体动作值函数 $Q_i(o_i,a_i)$}（值分解文献里常称 utility 效用，以区别于真正的 $Q$ 值），执行时各自选 $a_i^*=\arg\max_{a_i}Q_i(o_i,a_i)$。这样决策分布式（每人只在自己的 $|\mathcal{A}_i|$ 个动作里搜，线性复杂度）。但马上有个尖锐问题：\textbf{这些独立选出来的 $a_i^*$ 拼成的联合动作，凭什么是团队最优的}？我们真正想最大化的团队回报对应某个联合值 $Q_{\text{tot}}(s,\mathbf{a})$，希望 $(\arg\max_{a_1}Q_1,\dots,\arg\max_{a_N}Q_N)=\arg\max_{\mathbf{a}}Q_{\text{tot}}(s,\mathbf{a})$——但这个等式\emph{不会自动成立}，它对 $Q_i$ 和 $Q_{\text{tot}}$ 之间的关系提出了一个很强的结构要求。值分解这条线的全部技术内容，就是回答"$Q_i$ 与 $Q_{\text{tot}}$ 满足什么结构，这个等式才成立、且 $Q_{\text{tot}}$ 还能有足够表达力"。

\paragraph{反面：不要求这个等式 $\Rightarrow$ 局部理性、全局愚蠢。}
若随便学 $N$ 个 $Q_i$、不管它们和团队目标的关系，就退化成独立 Q-learning——局部最优拼出来的联合动作可能是团队灾难。举例：两智能体团队奖励"同 L 或同 R 得 $+10$、一 L 一 R 得 $0$"，若一个学会偏好 L、一个偏好 R，执行时各自贪心，拼出"一 L 一 R"的零回报灾难——明明各自都"选了自己认为最好的"，团队却最差。值分解要做的，就是给 $Q_i$ 套上结构，让局部贪心自动导向全局协调。

\subsection{VDN：最简的可加分解}

VDN（Value-Decomposition Networks）\cite{sunehag2018vdn}给出最简单的答案：假设联合值就是个体 utility 的\textbf{直接相加}：
\begin{equation}
  Q_{\text{tot}}(s,\mathbf{a})=\sum_{i=1}^N Q_i(o_i,a_i).
  \label{eq:sm-vdn}
\end{equation}
训练时把这个求和当成联合 $Q$，用标准（集中式）TD 误差以团队奖励 $r$ 为信号训练：$\mathcal{L}(\theta)=\mathbb{E}[(r+\gamma\max_{\mathbf{a}'}Q_{\text{tot}}(s',\mathbf{a}';\theta^-)-Q_{\text{tot}}(s,\mathbf{a};\theta))^2]$。注意梯度怎么流：$Q_{\text{tot}}=\sum_i Q_i$，TD 误差对 $Q_{\text{tot}}$ 的梯度通过求和\emph{自动分配}到每个 $Q_i$——这就是隐式的信用分配（\cref{sec:sm-credit} 会回到这点）。

\textbf{为什么可加分解能保证"局部 argmax $=$ 全局 argmax"？} 因为求和是\emph{可分离的}——每个 $a_i$ 只出现在自己那一项 $Q_i$ 里，不同智能体的动作在 $Q_{\text{tot}}$ 里没有交叉项。所以最大化整个和，等价于独立最大化每一项：$\arg\max_{\mathbf{a}}\sum_i Q_i(o_i,a_i)=(\arg\max_{a_1}Q_1,\dots,\arg\max_{a_N}Q_N)$。\textbf{VDN 的局限}：可加性太强——$Q_{\text{tot}}=\sum_i Q_i$ 意味着团队价值是个体价值的简单线性叠加，\emph{没有任何交互项}。但很多协作任务的价值本质上非线性耦合（如"两个智能体必须同时到达某点才触发奖励"——一个到了另一个没到，价值是 0 不是"半个奖励"）。而且 VDN 完全没用上全局状态 $s$。这两个局限催生了 QMIX。

\subsection{IGM 原则：把"等式"抽象成统一条件}

在讲 QMIX 前，先把 VDN 满足的那个"局部 argmax $=$ 全局 argmax"性质\textbf{抽象成一般原则}——这是值分解理论的核心，叫 IGM。

\begin{definition}[IGM 原则（Individual-Global-Max，个体-全局最大一致性）]\label{def:sm-igm}
一组个体 utility $\{Q_i(o_i,a_i)\}_{i=1}^N$ 与联合值 $Q_{\text{tot}}(s,\mathbf{a})$ 满足 IGM，当且仅当对所有状态：
\begin{equation}
  \arg\max_{\mathbf{a}}\ Q_{\text{tot}}(s,\mathbf{a})=\Big(\arg\max_{a_1}Q_1(o_1,a_1),\ \dots,\ \arg\max_{a_N}Q_N(o_N,a_N)\Big).
  \label{eq:sm-igm}
\end{equation}
即联合值的全局最优联合动作，恰好等于各个体 utility 各自局部最优动作的拼接。
\end{definition}

IGM 是值分解方法\emph{正确性的充要条件}——只要满足 IGM，执行时分布式地各自贪心就保证选出全局最优联合动作；不满足，分布式执行就可能选出次优甚至灾难性的联合动作。所以值分解方法的设计目标就变成了：\textbf{找一个 $Q_{\text{tot}}$ 与 $\{Q_i\}$ 的函数关系，使 IGM 成立、同时 $Q_{\text{tot}}$ 的表达力尽可能大}。VDN 用"求和"满足 IGM（可分离 $\Rightarrow$ IGM），但求和只是充分条件之一、表达力受限。

\begin{insight}
IGM 是值分解的"宪法"——它定义了什么样的分解是"合法的"（分布式执行能选出全局最优）。VDN 和 QMIX 不是两个孤立的算法，而是 IGM 这部宪法下的两种"实现方案"：VDN 用可加性、QMIX 用单调性来满足宪法。理解了 IGM，你就理解了整条值分解线——后面的 QTRAN、QPLEX 也都是在回答同一个问题"如何在满足 IGM 的前提下最大化表达力"，只是答案越来越精巧。
\end{insight}

\subsection{QMIX：单调混合网络}

QMIX\cite{rashid2018qmix}的核心洞察是：\textbf{IGM 只关心 $\arg\max$，不关心 $Q_{\text{tot}}$ 的具体数值形状}。而一个比"可加"宽得多、仍能保证 $\arg\max$ 可分解的条件是——$Q_{\text{tot}}$ 对每个 $Q_i$ \textbf{单调递增}：
\begin{equation}
  \frac{\partial\,Q_{\text{tot}}}{\partial\,Q_i}\ \ge\ 0,\quad\forall i\in\{1,\dots,N\}.
  \label{eq:sm-qmix-mono}
\end{equation}

\begin{theorem}[单调性蕴含 IGM]\label{thm:sm-mono-igm}
若 $Q_{\text{tot}}$ 对每个 $Q_i$ 单调递增（\cref{eq:sm-qmix-mono}），则 $\{Q_i\}$ 与 $Q_{\text{tot}}$ 满足 IGM（\cref{eq:sm-igm}）。
\end{theorem}
\begin{proof}
设 $a_i^*=\arg\max_{a_i}Q_i(o_i,a_i)$，记 $\mathbf{a}^*=(a_1^*,\dots,a_N^*)$。对任意联合动作 $\mathbf{a}=(a_1,\dots,a_N)$，因为 $Q_i(o_i,a_i^*)\ge Q_i(o_i,a_i)$ 对每个 $i$ 成立，又由 $Q_{\text{tot}}$ 对每个 $Q_i$ 单调递增，把 $\mathbf{a}$ 的分量逐个替换成 $a_i^*$（每替换一个，对应的 $Q_i$ 不减、由单调性 $Q_{\text{tot}}$ 不减）都不会减小 $Q_{\text{tot}}$，故 $Q_{\text{tot}}(s,\mathbf{a}^*)\ge Q_{\text{tot}}(s,\mathbf{a})$。即 $\mathbf{a}^*$ 最大化 $Q_{\text{tot}}$，IGM 成立。
\end{proof}

这个证明清晰地显示：\textbf{单调性是比可加性弱得多的要求}（可加是单调的特例：$\partial(\sum Q_j)/\partial Q_i=1\ge0$），却足以保证 IGM。所以 QMIX 的表达力严格大于 VDN，又不丢失分布式执行的正确性。

\paragraph{怎么参数化一个"对每个 $Q_i$ 单调递增、还能依赖全局状态 $s$"的混合网络？} 这是 QMIX 的工程精髓，两个设计。\textbf{设计 1：非负权重的混合网络（mixing network）}——用一个前馈网络把 $(Q_1,\dots,Q_N)$ 映射到 $Q_{\text{tot}}$；要保证对每个输入单调递增，只需让网络的\emph{所有权重非负}（外加单调递增激活如 ELU/绝对值），因为复合函数链式法则下每个偏导都是非负数的乘积与和，恒非负。\textbf{设计 2：用超网络（hypernetwork）注入全局状态 $s$}——这里有个矛盾：想让 $Q_{\text{tot}}$ 依赖 $s$（VDN 的缺陷正是用不上 $s$），但若把 $s$ 直接当混合网络的普通输入，要同时保证对 $Q_i$ 单调、对 $s$ 不必单调，权重约束会很别扭。QMIX 的巧妙解法是：\textbf{混合网络的权重不是固定参数，而是由另一个网络（超网络）以 $s$ 为输入生成的}。

\begin{codebox}[text]{QMIX 混合网络的数据流（超网络生成非负权重）}
   global state s
       |
   [hypernetwork]    arbitrary params; outputs pass |.| or ReLU -> nonneg
       |
       | generates mixing-net weights W1,W2 (nonneg) and biases b1,b2 (free)
       v
Q1,...,QN --> [mixing net: nonneg weights + monotone activation] --> Q_tot
              (weights come from hypernet, hence depend on s)
\end{codebox}

关键分工：\emph{混合网络}的权重 $W$ 被约束为\textbf{非负}（超网络输出过绝对值或 ReLU），保证 $\partial Q_{\text{tot}}/\partial Q_i\ge0$（IGM）；\emph{混合网络的偏置 $b$} 不需非负（偏置不影响对 $Q_i$ 的单调性，$\partial(Wx+b)/\partial x=W$ 与 $b$ 无关），可由超网络自由生成，让全局状态 $s$ 充分、非线性地影响 $Q_{\text{tot}}$ 的整体水平；\emph{超网络本身}参数任意（不受非负约束），它的输入是 $s$、输出是混合网络的权重——对 $s$ 怎么依赖都行，只要吐出的权重非负即可。这个设计精妙地分离了两个需求：\textbf{对 $Q_i$ 要单调（靠非负权重）、对 $s$ 要灵活（靠超网络）}。

\begin{note}
\textbf{一行代码承载一个定理.}\;QMIX 的非负权重约束在代码里就是一行 \texttt{torch.abs(w)} 或 \texttt{F.relu(w)}——把超网络输出的权重取绝对值/ReLU 再用。看起来微不足道，背后是整个 IGM 正确性的保证：少了这行，$Q_{\text{tot}}$ 不再对 $Q_i$ 单调，分布式执行不再选出全局最优，训练可能看起来在降 loss 但执行时智能体协调不起来。注意：只有\emph{权重}需要非负，\emph{偏置}不需要。
\end{note}

\subsection{QMIX 的局限与 QTRAN/QPLEX 的放松}

QMIX 用单调性满足 IGM，表达力比 VDN 大很多，但单调性仍是\emph{充分而非必要}条件——它能覆盖的 $Q_{\text{tot}}$ 函数类仍有缺口，经典反例是\textbf{非单调协调任务}。考虑两智能体单步博弈，团队奖励矩阵 $\big[\begin{smallmatrix}+8&-12\\-12&0\end{smallmatrix}\big]$（行 $a_1\in\{\text{A,B}\}$、列 $a_2\in\{\text{A,B}\}$）：最优联合动作是 $(\text{A,A})$ 回报 $+8$，但当 $a_2=\text{B}$ 时智能体 1 选 A 反而更差（$-12<0$）——即"$a_1$ 选 A 好不好"取决于 $a_2$ 选什么，这是\emph{非单调耦合}。QMIX 的单调性要求 $Q_{\text{tot}}$ 对 $Q_1$ 单调递增，迫使"$Q_1(\text{A})$ 大 $\Rightarrow Q_{\text{tot}}$ 大"对所有 $a_2$ 一致成立，但矩阵里 $a_1=\text{A}$ 的好坏随 $a_2$ 翻转，单调性结构无法拟合这种翻转。\textbf{这是单调性约束的表达力代价。}

\textbf{QTRAN}\cite{son2019qtran}的思路是：IGM 只要求 $\arg\max$ 一致，不要求 $Q_{\text{tot}}$ 处处等于某个单调函数。于是它学一个\emph{无结构约束}的联合 $Q_{\text{tot}}$（能表达任意耦合），同时学一组 $\{Q_i\}$，再通过\emph{两个线性约束}把它们"对齐"，理论上覆盖比 QMIX 更广的任务（含上面的非单调例子）。代价是那两个约束实践中难精确满足（要用罚项软化），训练更不稳、调参更难——所以 QTRAN 理论漂亮但工程上常打不过 QMIX。\textbf{QPLEX}\cite{wang2021qplex}用 dueling 架构（把 $Q$ 拆成 $V+A$，优势 $A$ 满足 $\max_a A=0$）重新表述 IGM，证明其分解能\emph{完整覆盖整个 IGM 函数类}（充要，而非 QMIX/VDN 的充分子集），同时保持训练稳定性——它是目前值分解表达力与稳定性平衡得最好的方法之一。

\begin{table}[H]\centering\small
\caption{值分解谱系：表达力 vs 可训练性}
\label{tab:sm-valuedecomp}
\begin{tabular}{p{0.10\textwidth} p{0.30\textwidth} p{0.22\textwidth} p{0.10\textwidth} p{0.10\textwidth}}
\toprule
方法 & $Q_{\text{tot}}$ 与 $\{Q_i\}$ 的关系 & 满足 IGM 的方式 & 表达力 & 可训练性 \\
\midrule
VDN & $Q_{\text{tot}}=\sum_i Q_i$ & 可加（可分离 $\Rightarrow$ IGM 充分） & 最弱 & 最稳 \\
QMIX & $Q_{\text{tot}}=f_{\text{mono}}(Q_1,\dots,Q_N;s)$ & 单调（IGM 充分，严格含可加） & 中 & 稳 \\
QTRAN & 无结构 $Q_{\text{tot}}$ $+$ 两个线性约束对齐 & 直接逼近 IGM（充要，靠软约束） & 强 & 不稳 \\
QPLEX & dueling 重述 $+$ 优势分解 & IGM 的完整充要刻画 & 强 & 较稳 \\
\bottomrule
\end{tabular}
\end{table}

\begin{insight}
从 VDN 到 QPLEX，这条线的内在逻辑是一个\textbf{不断逼近 IGM 边界}的过程。VDN 用很保守的充分条件（可加），留下大量 IGM 合法但 VDN 表达不了的分解；QMIX 把充分条件放宽到单调，覆盖更多但仍有缺口（非单调任务）；QTRAN/QPLEX 直接瞄准 IGM 的充要刻画。理解这一点，你就不会把这四个方法当四个孤立的"招式"记，而是看成同一个优化问题——"在 IGM 约束的可行域里最大化表达力"——的逐步精确求解。
\end{insight}

\begin{note}
\textbf{值分解线的工程定位.}\;\emph{适用}：离散动作（$\arg\max$ 是离散动作上的）、协作（共享奖励）、奖励相对稠密（稀疏奖励下 VDN/QMIX 学不动）。SMAC（星际争霸多智能体挑战）是经典主战场。\emph{不适用}：连续动作（无法直接 $\arg\max$，改用策略梯度）、奖励极稀疏、强异构需差异化角色。\textbf{机器人领域的现实}：大多数机器人控制任务是\emph{连续动作}（关节力矩、速度指令），这天然把值分解排除——这也是为什么 \cref{sec:sm-policygrad} 的策略梯度路线（尤其 MAPPO）才是机器人 MARL 的事实标准。值分解在机器人领域更多出现在\emph{离散高层决策}（子任务选择、运动基元切换）而非底层连续控制。
\end{note}

\begin{pitfall}[编程陷阱与概念误区：QMIX 的两个坑]
其一（编程）：\textbf{忘记给混合网络权重加非负约束}——超网络直接输出权重不经绝对值/ReLU 就用。现象：训练 loss 看起来在降（联合 TD 能拟合），但执行时智能体协调不起来、回报上不去（IGM 被破坏），这种"训练指标好、执行结果差"的脱节极难 debug。修复：超网络输出权重后\emph{必须}过 \texttt{torch.abs} 或 \texttt{F.relu}（偏置不需要），可在单元测试里数值检查 $\partial Q_{\text{tot}}/\partial Q_i$ 的符号。其二（概念）：\textbf{以为 QMIX 表达力够强能拟合任意协作任务}——在非单调协调任务（如上述 $\pm12$ 矩阵）上 QMIX 系统性把最优联合动作学偏，却以为是超参没调好。识别任务是否有非单调耦合（某智能体动作好坏强烈依赖他人且方向翻转），若有，换 QTRAN/QPLEX 或策略梯度。其三（思维）：把个体 utility $Q_i$ 当智能体 $i$ 的"真实价值"解读或复用——$Q_i$ 是 utility 不是 $Q$ 值，它没有独立语义（不满足任何单智能体贝尔曼方程），只在与混合网络组合后整体满足联合贝尔曼方程时才有意义，绝对尺度可被任意重标定；只用它做 $\arg\max_{a_i}Q_i$ 选动作，要评估单智能体贡献用专门归因（反事实/Shapley，\cref{sec:sm-credit}）。
\end{pitfall}

\begin{practice}
\begin{enumerate}
  \item (a) 证明可加分解是单调分解的特例（可加 $\Rightarrow$ 单调）；(b) 构造一个\emph{单调但不可加}的 $Q_{\text{tot}}(Q_1,Q_2;s)$ 具体例子并验证它满足 IGM；(c) 用 $\pm12$ 非单调矩阵论证不存在任何单调 $f$ 使 $Q_{\text{tot}}=f(Q_1,Q_2)$ 能正确表达该矩阵的 $\arg\max$ 结构。
  \item 用 PyTorch 实现最小 QMIX 混合网络：超网络生成两层混合网络的权重（过绝对值保非负）和偏置（不约束），输出 $Q_{\text{tot}}$；写数值测试用有限差分验证 $\partial Q_{\text{tot}}/\partial Q_i\ge0$，再故意去掉绝对值约束观察测试如何失败。
  \item 对比 CBBA 的"边际收益"分配与值分解的"梯度分摊"信用归因：(a) 思想上的相通（都是某种"边际贡献"）；(b) 关键区别（离散静态分配的显式边际 vs 序贯学习的隐式梯度信用）。
\end{enumerate}
\end{practice}

% ════════════════════════════════════════════════════════════════
\section{多智能体策略梯度：MADDPG → COMA → MAPPO → HAPPO}
\label{sec:sm-policygrad}

CTDE 骨架的第二种填法——集中式 critic $+$ 分布式 actor。这条路线是机器人 MARL 的事实标准（因为支持连续动作）。沿四步演进讲透：MADDPG 用集中 critic 吸收非平稳性（解困难①），COMA 用反事实基线解信用分配（解困难②），MAPPO 用"共享 PPO $+$ 全局 critic"成为最强基线（解①③），HAPPO 用多智能体优势分解引理实现异构智能体的单调改进（解④）。每一步都建立在单智能体策略梯度定理之上——这条线就是策略梯度定理的"多智能体加料版"。

\paragraph{动机：值分解的天花板与策略梯度的必要性。}
值分解很优美，但有个硬伤——它要 $\arg\max_{a_i}Q_i$，只在\emph{离散动作}上可行。机器人控制几乎都是\emph{连续动作}（关节力矩 $\tau\in\mathbb{R}^d$、速度指令），在连续空间上 $\arg\max$ 一个 $Q$ 本身就是非凸优化、做不动。所以机器人 MARL 必须走\textbf{策略梯度}——直接参数化策略 $\pi_i(a_i\mid o_i)$ 并对其参数做梯度上升，绕开 $\arg\max$。回顾单智能体策略梯度定理 $\nabla_\theta J(\theta)=\mathbb{E}_{s\sim d^\pi,a\sim\pi_\theta}[\nabla_\theta\log\pi_\theta(a\mid s)\,A^\pi(s,a)]$。最朴素的多智能体写法是给每个 $i$ 套一份单智能体公式（这就是 IPPO/独立 actor-critic）：$\nabla_{\theta_i}J\approx\mathbb{E}[\nabla_{\theta_i}\log\pi_{\theta_i}(a_i\mid o_i)\,A_i]$。但这里有两个空白要填：(1) $A_i$ 怎么算——用什么 critic、用团队奖励还是个体贡献？(2) 直接独立做会撞上非平稳性。四步演进就是逐步把这两个空白填好。

\subsection{第一步：MADDPG——集中式 critic 吸收非平稳性}

MADDPG\cite{lowe2017maddpg}是 CTDE 的奠基作。它从连续动作的确定性策略梯度（DDPG）出发，核心创新是 critic 的集中化。\textbf{架构}：每个智能体 $i$ 有一个\emph{分布式 actor} $\mu_{\theta_i}(o_i)$（确定性策略，吃本地观测出连续动作，执行时只用它），和一个\emph{集中式 critic} $Q_i^\mu(s,a_1,\dots,a_N)$（吃全局状态 $s$ 和所有人的动作，只在训练时用）。\textbf{为什么 critic 要吃所有人的动作？} 这正是 \cref{sec:sm-ctde} 推过的——把别人的动作 $\mathbf{a}_{-i}$ 显式作为条件，非平稳性被吸收：对固定的 $(s,\mathbf{a})$，转移核真正固定，所以 $Q_i(s,\mathbf{a})$ 满足良定义的贝尔曼方程、其 TD 学习稳定。注意 MADDPG 给每个 $i$ 一个\emph{独立}的 critic $Q_i$（而非共享）——这是为了支持\emph{一般和}博弈（每个智能体可有不同奖励 $r_i$）；纯合作下所有 $Q_i$ 会趋同，可共享一个。\textbf{梯度（确定性策略梯度版）}：
\begin{equation}
  \nabla_{\theta_i}J_i=\mathbb{E}_{s,\mathbf{a}\sim\mathcal{D}}\Big[\nabla_{\theta_i}\mu_{\theta_i}(o_i)\,\nabla_{a_i}Q_i^\mu(s,a_1,\dots,a_N)\big|_{a_i=\mu_{\theta_i}(o_i)}\Big].
  \label{eq:sm-maddpg}
\end{equation}
直觉：固定别人的动作（从 replay buffer $\mathcal{D}$ 里取），问"我把自己的动作 $a_i$ 往哪个方向微调，能让集中 critic 的 $Q_i$ 升高"，让 actor 朝那个方向更新；critic 通过标准 DDPG 的 TD 误差训练，因为它吃联合动作，旧 replay 数据不过期。

\begin{insight}
MADDPG 的全部精华在一句话——\textbf{"actor 局部、critic 全局"}。actor 必须局部（执行约束），critic 可以全局（训练时仿真器给得起全局信息）。这个不对称把非平稳性从 actor 的学习里剥离出来、塞进了 critic 的输入条件。MADDPG 之后所有 CTDE actor-critic 方法（COMA、MAPPO、HAPPO）都继承这个核心，只在"critic 用什么形式、优势怎么算、怎么更新"上分化。
\end{insight}

\subsection{第二步：COMA——反事实基线解信用分配}

COMA\cite{foerster2018coma}专攻信用分配（困难②）。它是 on-policy actor-critic，用一个集中式 critic，关键创新在\textbf{基线（baseline）的选择}。合作场景只有团队奖励 $r$，若 $A_i$ 用团队优势（对所有智能体一样），$i$ 的梯度被"别人的贡献"污染。回顾单智能体基线的作用——减去一个不依赖动作 $a_i$ 的项，降方差且不引入偏差。COMA 的天才之处是选一个特殊的、依赖全局信息的基线——把智能体 $i$ 的动作边缘化掉、固定其他人动作时的期望 $Q$：
\begin{equation}
  A_i(s,\mathbf{a})=Q(s,\mathbf{a})-\underbrace{\sum_{a_i'}\pi_i(a_i'\mid\tau_i)\,Q\big(s,(a_i',\mathbf{a}_{-i})\big)}_{\text{反事实基线：固定别人、平均掉 }i\text{ 自己的动作}}.
  \label{eq:sm-coma}
\end{equation}
逐项拆解：$Q(s,\mathbf{a})$ 是当前联合动作的联合值；基线项保持别人动作 $\mathbf{a}_{-i}$ 不变、把 $i$ 的动作换成它的策略平均；两者之差 $A_i$ 回答反事实问题——\textbf{"$i$ 实际选的这个 $a_i$，比 $i$ 平均水平的动作好多少？"} 这个差值剔除了别人的贡献（别人的动作在被减的两项里都一样、抵消了），纯粹反映 $i$ 自己这步动作相对自己基准的优劣。\textbf{为什么不引入偏差？} 因为基线项不依赖 $i$ 实际采的动作 $a_i$（对 $i$ 是个"常数基线"），由策略梯度的基线定理，减去它不改变梯度期望、只降方差。COMA 用一个集中 critic 一次前向就给出 $i$ 所有动作的 $Q$，故反事实基线高效。

\begin{note}
\textbf{COMA $\neq$ 平分 / 手工塑形.}\;反事实基线不是"把团队奖励除以 $N$ 平分"，也不是手工给每个智能体设计局部奖励。它是用集中 critic 估计一个\emph{反事实量}——"如果 $i$ 换个动作、别人不变，团队会差多少"。这是差分奖励（\cref{sec:sm-credit}）的思想用 critic 实现的版本：差分奖励要重置仿真器算反事实（昂贵、需可重置仿真器），COMA 用一次 critic 前向就近似出反事实。\textbf{关键边界}：COMA 的反事实基线只对离散动作直接可算（要对 $a_i'$ 求和），连续动作要积分、难做——这是 COMA 主要用于离散动作的原因。实践中（MAPPO 基准）COMA 常打不过更简单的 MAPPO——反事实 critic 难训、方差仍大。它的价值更多在\emph{思想}（反事实基线是信用分配的范式），而非作为首选算法。
\end{note}

\subsection{第三步：MAPPO——共享 PPO \texorpdfstring{$+$}{+} 全局 critic，最强基线}

MAPPO\cite{yu2022mappo}是当前合作 MARL（尤其机器人领域）的事实标准基线。它的故事有点"反高潮"——几乎没有新理论，就是把单智能体 PPO 几乎原封不动搬到多智能体、加两个朴素设计，结果在 SMAC、Multi-Agent MuJoCo、Hanabi 上\emph{全面胜过}精心设计的 QMIX、MADDPG、COMA。

\textbf{设计 1：全局 critic.} critic 是一个\emph{状态值函数} $V_\phi(s)$，吃全局状态 $s$（CTDE 的集中信息），输出标量值——注意是 $V(s)$ 而非 $Q(s,\mathbf{a})$，它不吃动作。优势用广义优势估计（GAE）算：
\begin{equation}
  A_i^{\text{GAE}}=\sum_{l=0}^{\infty}(\gamma\lambda)^l\,\delta_{t+l},\qquad \delta_t=r_t+\gamma V_\phi(s_{t+1})-V_\phi(s_t).
  \label{eq:sm-mappo-gae}
\end{equation}
\textbf{设计 2：参数共享.} 所有智能体共享同一个 actor 网络 $\pi_\theta$ 和同一个 critic 网络 $V_\phi$，actor 吃本地观测 $\pi_\theta(a_i\mid o_i)$；因为参数共享，所有智能体的经验都用来更新同一套参数——\emph{样本效率 $N$ 倍提升}。actor 目标就是 PPO 的 clip（逐智能体）：
\begin{equation}
  L^{\text{actor}}(\theta)=\mathbb{E}_i\Big[\min\big(\rho_i A_i,\ \mathrm{clip}(\rho_i,1-\epsilon,1+\epsilon)A_i\big)\Big],\quad \rho_i=\frac{\pi_\theta(a_i\mid o_i)}{\pi_{\theta_{\text{old}}}(a_i\mid o_i)}.
  \label{eq:sm-mappo-clip}
\end{equation}
这和单智能体 PPO 的 clip 目标\emph{逐字一样}，只是期望多了对智能体 $i$ 的平均。\textbf{整个 MAPPO 训练循环 $=$ 单智能体 PPO 循环 $+$ "用全局 state 喂 critic" $+$ "把 $N$ 个智能体的数据合并训练"}。

\textbf{为什么这么简单的方法这么强？} 几个原因：(1) \emph{on-policy 避开数据过期}——PPO 每次用最新策略采的数据、用完即弃，不存 replay buffer，从源头规避多智能体 off-policy 的数据过期（\cref{sec:sm-nonstationary} 破坏二），这是它比 off-policy 的 MADDPG/QMIX 更稳的重要原因；(2) \emph{全局 critic 治非平稳 $+$ 降方差}；(3) \emph{参数共享提样本效率}；(4) \emph{PPO 的 clip 本身就鲁棒}——信赖域裁剪让更新保守、不容易因为一次坏的优势估计崩掉，这种鲁棒性在更不稳定的多智能体环境里更值钱。

\begin{insight}
MAPPO 的"惊人有效"传递了一个深刻的方法论信号——\textbf{在多智能体里，简单 $+$ 鲁棒常常胜过精巧 $+$ 脆弱}。QMIX/COMA 用精巧结构（单调混合、反事实基线）针对性地解某个困难，但这些结构本身引入新的训练不稳定性和调参负担；MAPPO 用最朴素的"共享 PPO $+$ 全局 critic"，每个部件都久经单智能体考验、极其鲁棒，组合起来反而最稳最强。这不是说精巧方法没价值（它们贡献思想、在特定任务上更优），而是说\textbf{当你不确定用什么时，MAPPO 是正确的默认起点}——尤其在机器人连续控制上。
\end{insight}

\begin{note}
\textbf{决定 MAPPO 成败的关键实现技巧.}\;\cite{yu2022mappo} 的贡献一半在算法、一半在系统性消融出几个关键技巧：\emph{value normalization（值归一化）}——用 running mean/std 归一化 critic 的回报目标，对不同量级奖励至关重要；\emph{全局 state 的构造}——用真正的全局状态（agent-specific global state 常更好），而非简单 concat 所有观测（\cref{sec:sm-marl-form} 陷阱）；PPO clip 系数、GAE $\lambda$、并行环境数、训练 epoch 数与 mini-batch（多智能体下 over-fitting 一个 batch 的风险更高，epoch 数通常比单智能体小）。\textbf{MAPPO 的边界}：它依赖参数共享 $\Rightarrow$ 要求同构（所有智能体观测/动作空间相同、角色可互换）。异构智能体（不同传感器、不同动作空间、不可互换角色）下参数共享不适用——这正是 HAPPO 要解决的。
\end{note}

\subsection{第四步：HAPPO——多智能体优势分解引理与异构单调改进}

HAPPO\cite{kuba2022happo}解决 MAPPO 的两个遗留问题：(a) 参数共享要求同构，异构智能体怎么办？(b) MAPPO 没有单调改进的理论保证（PPO 单智能体有近似单调改进保证，但多智能体同时更新所有 actor 时这个保证失效）。HAPPO 的核心是一个漂亮的引理。

\textbf{问题：为什么不能所有 actor 同时更新？} 单智能体 PPO 的单调改进保证（来自 TRPO）说：在信赖域内更新一个策略，性能单调不降。但多智能体里，若所有 $N$ 个 actor 同时各自在自己的信赖域内更新，\emph{它们的联合变化可能跳出联合信赖域}——每个人单独看是小步，合起来是大步，单调改进保证失效（这其实又是非平稳性的一个化身：你更新时假设别人不变，但别人也在同时变）。

\begin{theorem}[多智能体优势分解引理]\label{thm:sm-advantage-decomp}
先定义多智能体优势：对智能体子集，给定前面智能体的动作，后一个智能体的优势 $A^{i_{1:k}}(s,a^{i_{1:k}}):=Q^{i_{1:k}}(s,a^{i_{1:k}})-Q^{i_{1:k-1}}(s,a^{i_{1:k-1}})$，其中 $Q^{i_{1:k}}$ 是固定前 $k$ 个智能体动作、其余按当前策略平均的联合 $Q$。则对任意智能体排列 $i_1,\dots,i_N$，联合优势可\emph{精确}分解为
\begin{equation}
  A(s,a_1,\dots,a_N)=\sum_{k=1}^{N}A^{i_{1:k}}\big(s,\ a^{i_{1:k}}\big).
  \label{eq:sm-adv-decomp}
\end{equation}
\end{theorem}

\textbf{这个引理在说什么？} 它说联合优势可以\emph{无损地}写成一个"逐个加入智能体"的累积：第一个智能体贡献 $A^{i_1}$（它单独行动相对平均的优势），第二个贡献 $A^{i_1 i_2}$（在第一个动作已定的前提下、它的边际优势），以此类推。这是一个\emph{精确等式}（不是近似），且对任意更新顺序成立。其结构和概率链式法则 $P(x_1,\dots,x_N)=\prod_k P(x_k\mid x_{1:k-1})$ 完全同构——只是把乘积换成求和、把概率换成优势。

\textbf{HAPPO 的顺序更新.} 有了分解引理，HAPPO 不再同时更新所有 actor，而是\emph{按一个随机顺序逐个更新}：随机排列智能体顺序 $i_1,\dots,i_N$；更新第一个 $i_1$（用条件优势 $A^{i_1}$ 做 PPO 更新），固定其余；更新第二个 $i_2$，此时 $i_1$ 已是新策略（把 $i_1$ 的更新"告诉" $i_2$），用条件优势 $A^{i_1 i_2}$ 更新；依次到 $i_N$。\textbf{为什么顺序更新能保证单调改进？} 因为每一步只更新一个智能体、固定其余——这回到了"单智能体在固定环境里更新"，单智能体 PPO/TRPO 的单调改进保证适用：这一步不会降低联合性能。而分解引理保证把这些"逐个的改进"加起来恰好等于联合优势的改进。于是一整轮（所有智能体各更新一次）下来，联合性能\emph{单调不降}——这是 MAPPO 没有的理论保证。其逻辑结构和坐标下降一模一样：一次只优化一个坐标（智能体）、固定其余，每步不降、整体收敛。

\begin{insight}
HAPPO 用"顺序"换"保证"。MAPPO 同时更新所有 actor（并行、快）但放弃单调改进保证（实践中靠 PPO 的 clip 凑合）；HAPPO 顺序更新（串行、慢一点）换来严格的单调改进保证 $+$ 天然的异构支持（每个智能体独立 actor/critic，不要求同构）。坐标下降的智慧在这里再次显现：把一个难的联合优化，拆成一串容易的、有保证的单变量优化。
\end{insight}

\begin{table}[H]\centering\small
\caption{MAPPO vs HAPPO 选型}
\label{tab:sm-mappo-happo}
\begin{tabular}{p{0.16\textwidth} p{0.36\textwidth} p{0.36\textwidth}}
\toprule
维度 & MAPPO & HAPPO \\
\midrule
参数 & 共享（同一套 actor/critic） & 每个智能体独立 actor/critic \\
同构要求 & 要求同构（观测/动作空间相同） & 天然支持异构 \\
更新方式 & 所有 actor 同时更新（并行） & 顺序逐个更新（串行） \\
理论保证 & 无单调改进保证 & 有单调改进保证（分解引理） \\
样本效率 & 高（参数共享，$N$ 倍数据） & 较低（独立网络，数据不共享） \\
适用 & 同构协作（编队、同型机群） & 异构协作（领航-跟随、不同机型、人形不同肢体） \\
\bottomrule
\end{tabular}
\end{table}
一句话选型：\textbf{同构用 MAPPO（简单高效），异构用 HAPPO（有保证、支持差异化角色）}。HARL 库\cite{zhong2024harl}统一实现了 HAPPO/HATD3/HASAC 等异构方法，是异构 MARL 的标准代码库。

\begin{pitfall}[编程陷阱：HAPPO 的顺序更新不是"简单循环"]
实现 HAPPO 时简单地 \texttt{for i in agents: update(i)}、但每个智能体用的还是同一批一开始算好的优势、不重算重要性采样比，会\emph{失去单调改进保证}。根本原因：分解引理的顺序更新要求——更新智能体 $i_k$ 时，前面 $i_1,\dots,i_{k-1}$ 已变成新策略，$i_k$ 看到的"其余智能体联合策略"变了，所以 $i_k$ 的重要性采样比 $\rho_{i_k}$ 必须\emph{基于已更新的前序策略重新计算}。少了这个"逐步重算"就不是真正的 HAPPO，分解引理的保证不成立。正确做法：随机排列后逐个更新，每更新完一个，在更新下一个之前用最新的（含已更新前序的）联合策略重算后续智能体的重要性比和优势（HARL 库的 factor 累积逻辑）。另一个相关误区：把 MAPPO 的 critic 当成吃联合动作的 $Q(s,\mathbf{a})$——MAPPO 的 critic 是\emph{状态值函数 $V_\phi(s)$}（只吃全局状态、不吃动作），优势用 GAE，走的是 A2C/PPO 的"$V+$GAE"路线，不是 DDPG 的"$Q+$确定性梯度"路线。
\end{pitfall}

\subsection{两条路线总对比与选型}

把值分解和策略梯度两条 CTDE 路线做头对头对比，这是 MARL 选型时最常回查的一张表：

\begin{table}[H]\centering\small
\caption{值分解 vs 策略梯度两条路线}
\label{tab:sm-two-lines}
\begin{tabular}{p{0.18\textwidth} p{0.36\textwidth} p{0.36\textwidth}}
\toprule
维度 & 值分解（VDN/QMIX/QPLEX） & 策略梯度（MADDPG/MAPPO/HAPPO） \\
\midrule
集中组件 & 混合网络（聚合 $Q_i\to Q_{\text{tot}}$） & 集中 critic（$V(s)$ 或 $Q(s,\mathbf{a})$） \\
学的是什么 & 个体 utility $Q_i$（值） & 策略 $\pi_i$ $+$ critic（值） \\
动作空间 & 仅离散（依赖 $\arg\max_{a_i}$） & 离散 $+$ \textbf{连续}（机器人控制必需） \\
正确性依赖 & IGM 原则（单调性约束） & 策略梯度定理（无 IGM 约束） \\
信用分配 & 隐式（联合 TD 误差梯度分摊） & 显式（团队优势 / 反事实基线） \\
on/off-policy & off-policy（用 replay，要处理过期） & MAPPO/COMA on-policy；MADDPG off-policy \\
奖励稠密度 & 需稠密（稀疏下学不动） & 对稀疏奖励更鲁棒 \\
典型主战场 & SMAC、离散微操协作 & 机器人连续控制、Multi-Agent MuJoCo \\
"最强"方法 & QMIX（稠密离散）/ QPLEX（表达力） & MAPPO（同构）/ HAPPO（异构） \\
\bottomrule
\end{tabular}
\end{table}

\begin{insight}
两条路线不是"谁取代谁"，而是\textbf{各占一个生态位}——值分解占"离散 $+$ 稠密 $+$ 需要值函数"的生态位（SMAC 是它的家），策略梯度占"连续 $+$ 机器人控制"的生态位（这是更大的、机器人最关心的生态位）。它们共享同一个 CTDE 骨架，区别只在"集中组件是混合网络还是 critic、学的是值还是策略"。理解这一点，你就不会问"QMIX 和 MAPPO 到底谁好"这种没有绝对答案的问题，而会问"我的任务落在哪个生态位"。对机器人工程师：\textbf{默认从 MAPPO 起步}，因为机器人控制几乎都是连续动作，值分解从第一刀就出局了。
\end{insight}

\begin{practice}
\begin{enumerate}
  \item 拿一份单智能体 PPO 实现，列出改成 MAPPO 的\emph{最小改动清单}：(a) critic 输入从 $s$ 改成什么？(b) 数据收集如何从 1 个智能体扩展到 $N$ 个？(c) 优势/loss 计算需要改吗？(d) actor 是否需要 agent ID？逐条说明每个改动对应解决 \cref{sec:sm-four-difficulties} 的哪个困难。
  \item 证明 COMA 的反事实基线满足"基线无偏性"——$\mathbb{E}_{a_i\sim\pi_i}[\nabla_{\theta_i}\log\pi_i(a_i\mid\tau_i)\,b_i(s,\mathbf{a}_{-i})]=0$（提示：$b_i$ 不依赖 $a_i$，套用 $\sum_{a_i}\pi_i\nabla\log\pi_i=\nabla\sum_{a_i}\pi_i=\nabla 1=0$）。
  \item 取 $N=2$ 合作单步博弈，自定 $Q(s,a_1,a_2)$ 和策略，按 \cref{thm:sm-advantage-decomp} 计算 $A^1$ 和 $A^{12}$（顺序 $1\to2$），验证 $A^1+A^{12}=A(s,a_1,a_2)$；换顺序 $2\to1$ 再算，确认联合优势相同但中间各项不同。
\end{enumerate}
\end{practice}

% ════════════════════════════════════════════════════════════════
\section{信用分配：团队奖励下"谁的功劳"}
\label{sec:sm-credit}

信用分配（credit assignment）在前面几节反复出现——VDN/QMIX 的梯度分摊、COMA 的反事实基线都在做它。这一节把它单独拎出来做\textbf{统一视角}：把信用分配形式化为一个反事实归因问题，揭示差分奖励、反事实基线、值分解三种解法其实是同一思想的三种实现，并引入 Shapley 值作为"理论上正确的归因"的标尺。

\paragraph{动机：搭便车问题。}
回到 \cref{sec:sm-four-difficulties} 困难②的场景：多机搬运只有一个团队奖励——货送到给 $+100$。三个机器人，1 号干了 90\% 的活、2 号干了 10\%、3 号全程摸鱼。团队成功了，团队奖励 $+100$。若三个都用同一个 $+100$ 当信号，那 3 号的"摸鱼"动作也被 $+100$ 强化了，下次 3 号会更倾向摸鱼——这就是\textbf{搭便车（free-riding）}。更微妙的反面也存在：\textbf{懒惰智能体问题（lazy agent，QMIX 论文专门讨论\cite{rashid2018qmix}）}——如果一个智能体一旦行动就容易把事情搞砸（高方差），团队可能学出"让它别动"的策略，于是它永远学不到有用行为、退化成废物。两个现象都是信用分配失败的症状：\textbf{团队奖励没有正确地反映每个智能体的真实边际贡献}。最朴素的两种办法都有致命缺陷：\emph{团队奖励平分}（每人 $r/N$）完全没区分贡献；\emph{手工设计个体奖励（reward shaping）}不可扩展、易引入偏差（诱发奖励黑客 reward hacking）、把本该自动解决的问题甩给人类。问题的关键在于：\textbf{贡献是一个反事实概念}——"$i$ 的贡献" $=$ "有 $i$ 参与时的团队结果" 减去 "没有 $i$（或 $i$ 换个动作）时的团队结果"。

\subsection{差分奖励：信用分配的"金标准"思想}

差分奖励（difference rewards）\cite{wolpert2002difference}给出了信用分配最干净的定义。智能体 $i$ 的差分奖励是：
\begin{equation}
  D_i(s,\mathbf{a})=r(s,\mathbf{a})-r\big(s,\,(c_i,\mathbf{a}_{-i})\big),
  \label{eq:sm-diff-reward}
\end{equation}
其中 $c_i$ 是智能体 $i$ 的一个"默认动作"（default action，比如"什么都不做"或某基准动作），$\mathbf{a}_{-i}$ 是其他人的实际动作。逐项理解：$r(s,\mathbf{a})$ 是所有人按实际动作时的团队奖励；$r(s,(c_i,\mathbf{a}_{-i}))$ 是把 $i$ 的动作换成默认 $c_i$、别人不变时的团队奖励（"$i$ 没真正参与时团队拿多少"）；两者之差 $D_i$ 是 $i$ 的实际动作相对默认动作给团队带来的增量——这就是 $i$ 的边际贡献。

\textbf{为什么差分奖励解决搭便车？} 摸鱼的 3 号，它的实际动作（摸鱼）和默认动作（什么都不做）几乎一样，所以 $D_3\approx0$——摸鱼拿到的信用接近零，不会被强化；而干活的 1 号，把它的动作换成默认会让团队结果大幅下降，所以 $D_1$ 很大。把搬运例子的数字代进去：默认动作 $c_i=$"什么都不做"，则 $D_1=100-0=100$（换 1 号没人扛主活、货送不到）、$D_2=100-90=10$（1 号扛 90\%，货大概率还能送到）、$D_3=100-100=0$（和它摸鱼没区别）。对比团队奖励 $100$ 平分三个机器人都拿 $33.3$，差分奖励给出 $(100,10,0)$——精确反映各自真实边际贡献，搭便车被结构性消除。

\begin{insight}
差分奖励揭示了信用分配的第一性原理——\textbf{贡献 $=$ 反事实差值}。"我的功劳"不是"团队总分的某个份额"，而是"有我和没我之间团队的差别"。这个反事实视角是整个信用分配领域的灵魂，后面 COMA 的反事实基线、值分解的梯度分摊都是它的不同实现。一旦你用反事实视角看信用分配，搭便车和懒惰智能体都迎刃而解——摸鱼者的反事实差值为零、被边缘化；真正有用者的反事实差值大、被强化。\textbf{差分奖励的工程难点}：它需要计算 $r(s,(c_i,\mathbf{a}_{-i}))$——要么重置仿真器重新跑一遍（昂贵、要求可重置到任意状态），要么有奖励函数的解析形式能直接代入。两者在真实任务里常常做不到，这正是 COMA 的反事实基线要绕开的。
\end{insight}

\subsection{三种解法的统一视角与 Shapley 标尺}

把前面散落的信用分配机制收拢，会发现它们都是"反事实差值"思想的不同实现，只是在"反事实怎么算"上分化：

\begin{table}[H]\centering\small
\caption{信用分配三种解法的统一视角}
\label{tab:sm-credit}
\begin{tabular}{p{0.16\textwidth} p{0.34\textwidth} p{0.20\textwidth} p{0.20\textwidth}}
\toprule
解法 & 反事实差值怎么估 & 优点 & 缺点 \\
\midrule
差分奖励 & 重置仿真器/解析奖励，算"$i$ 换默认动作"的团队奖励 & 概念最干净、归因最准 & 需可重置仿真器或解析奖励，常不可行 \\
反事实基线（COMA） & 用集中 critic $Q(s,\mathbf{a})$ 估"$i$ 换动作"的值，对 $a_i$ 求期望当基线 & 无需重置仿真器，一次 critic 前向 & critic 难训、离散动作为主 \\
值分解（VDN/QMIX） & 联合 TD 误差对 $Q_{\text{tot}}$ 求梯度，自动按 $\partial Q_{\text{tot}}/\partial Q_i$ 分摊 & 隐式、和值学习一体 & 受 IGM 约束、离散动作、吃稠密奖励 \\
\bottomrule
\end{tabular}
\end{table}

它们的内在统一：差分奖励直接、显式地算反事实差值（代价是要真的"执行反事实"）；COMA 的反事实基线用 critic \emph{预测}反事实而非真的执行（绕开仿真器重置，是差分奖励的"可微、可学习"版本，默认动作换成"策略平均动作"）；值分解更隐式——它不显式构造反事实，而是让团队 TD 误差的梯度\emph{自动}按每个 $Q_i$ 对 $Q_{\text{tot}}$ 的边际影响分摊（信用分配"免费"地内嵌在值学习的反向传播里，代价是受 IGM/单调性约束）。

\textbf{Shapley 值——归因的理论标尺.} 如果要问"什么是理论上唯一正确的贡献归因"，答案来自合作博弈论的 Shapley 值。它考虑智能体以\emph{所有可能的顺序}加入团队，$i$ 的 Shapley 值是它在所有顺序下边际贡献的平均：
\begin{equation}
  \phi_i=\sum_{\mathcal{C}\subseteq\mathcal{N}\setminus\{i\}}\frac{|\mathcal{C}|!\,(N-|\mathcal{C}|-1)!}{N!}\big[v(\mathcal{C}\cup\{i\})-v(\mathcal{C})\big],
  \label{eq:sm-shapley}
\end{equation}
其中 $v(\mathcal{C})$ 是子集 $\mathcal{C}$ 的智能体合作能达成的团队价值，方括号里是 $i$ 加入子集 $\mathcal{C}$ 的边际贡献。Shapley 值是\emph{唯一}同时满足"效率、对称、虚拟玩家、可加"四条公理的归因方案。差分奖励只用了\emph{一个}默认动作下的边际贡献（相当于只考虑一种"加入顺序/基准"），可看成 Shapley 值的粗糙近似；HAPPO 的优势分解引理（\cref{thm:sm-advantage-decomp}）按某个顺序逐个算边际优势，其结构也和 Shapley 的"按顺序加入算边际"同源——只是 HAPPO 用单一随机顺序、Shapley 用全顺序平均。\textbf{为什么实践中不直接用 Shapley 值？} 因为它要对 $2^N$ 个子集求和——指数复杂度，$N$ 稍大就算不动。所以实践中用各种近似，Shapley 值的价值在于给出"正确归因"的理论标尺。

\begin{note}
\textbf{与可解释 AI 的同源.}\;MARL 的信用分配之于团队，如同可解释 AI 里的特征归因之于模型预测。SHAP（Shapley Additive Explanations）用 Shapley 值归因"每个特征对预测贡献多少"，和 MARL 用 Shapley 值归因"每个智能体对团队回报贡献多少"是\emph{同一个数学工具的两个应用}。关键不同：可解释 AI 里参与者（特征）是被动的、归因只为解释；MARL 里参与者（智能体）是主动学习的、归因直接进入梯度去改变策略——归因错了不只是解释错，而是会把策略带偏（搭便车、懒惰智能体）。
\end{note}

\begin{pitfall}[思维陷阱：忽视"懒惰智能体"——信用分配失败的另一面]
只关注"摸鱼者被错误强化"（搭便车），忽视了反方向的失败——有用智能体因高方差被团队学会"让它别动"。现象：某个本该有用的智能体训练后退化成"什么都不做"，团队靠其余智能体硬撑。根本原因：当一个智能体一旦行动就容易（暂时）拉低团队回报（探索期高方差），团队级信用分配可能把"它行动"和"团队变差"错误关联，于是学出"抑制它"的策略，它便永远得不到正向信用、学不到有用行为。正确做法：(a) 用能正确归因边际贡献的信用分配（反事实/值分解），让有用智能体即使短期高方差也能获得正确的长期信用；(b) 保证充分探索（给每个智能体独立的探索噪声）；(c) 监控每个智能体的行为熵和贡献，及早发现退化为废物的智能体。另一个误区：以为给每个智能体单独奖励信号就解决了信用分配——手工局部奖励既不保证反映真实边际贡献、又把责任甩给人类，是"假装解决"。
\end{pitfall}

\begin{practice}
\begin{enumerate}
  \item 设计一个 3 智能体单步任务：团队奖励 $=$ 至少 2 个智能体选"合作"时 $+10$ 否则 $0$，每个选"合作"有 $-1$ 个人成本（不计入团队奖励，只影响真实意图）。(a) 论证三个都用团队奖励做策略梯度时为什么涌现搭便车；(b) 计算每个的差分奖励 $D_i$（默认动作 $=$"不合作"），说明差分奖励如何让"真正促成团队达标"的智能体获得正确信用。
  \item 对上题任务，分别写出：(a) 差分奖励 $D_i$ 的显式表达式；(b) COMA 反事实基线下 $i$ 的优势 $A_i$（假设有准确 $Q(s,\mathbf{a})$）；(c) 若用 VDN，团队 TD 误差如何通过 $Q_{\text{tot}}=\sum Q_i$ 分摊到每个 $Q_i$。对比三者得到的"信用"在什么假设下一致。
  \item 对一个 3 智能体合作博弈，给定所有 $2^3=8$ 个子集的价值 $v(\mathcal{C})$（自设一组数），手算每个智能体的 Shapley 值 $\phi_i$，验证效率公理 $\sum_i\phi_i=v(\{1,2,3\})$；再用"单一加入顺序的边际贡献"（差分奖励式近似）算一遍，对比它和 Shapley 值的差距。
\end{enumerate}
\end{practice}

% ════════════════════════════════════════════════════════════════
\section{MARL 三主线与 PSRO：从自博弈中学出均衡}
\label{sec:sm-psro}

前面两条 CTDE 路线（值分解、策略梯度）主攻\emph{合作}。但 MARL 与博弈论的交界更宽——还有\emph{竞争}与\emph{混合}场景。把 MARL 与博弈论交界的方法系统地分成三主线，并重点讲其中统一性最强的 PSRO。

\paragraph{动机：博弈太大，写不出也解不动。}
连续微分博弈的求解器（iLQGames 类）都假设你能\emph{显式写出}博弈：每个 agent 的代价函数、动力学、约束。少数 agent 的连续控制（自驾、竞速）里没问题，但有两类场景显式求解彻底失效：其一，\emph{博弈巨大到无法显式表示}——扑克的信息集有 $10^{160}$ 个，星际争霸的状态-动作空间天文数字，根本无法写出完整博弈树或代价函数；其二，\emph{代价函数本身未知或难手工设计}——你想训练一群无人机协作，但"什么是好的协作"很难写成精确代价。这两类场景需要换思路：\textbf{不显式求解博弈，而是用强化学习让 agent 在反复对弈（自博弈）中学出好策略，最终逼近博弈均衡}——AlphaGo、AlphaStar、Pluribus 都靠它在围棋、星际、扑克上达到超人水平。

\paragraph{反面：独立 RL 在多智能体里失效。}
最朴素的想法是独立 RL（Independent RL, InRL）——每个 agent 各自跑一个 RL、把其他 agent 当环境的一部分。它有两个根本问题：其一，\emph{环境非平稳}（\cref{sec:sm-nonstationary} 已详述）——你刚学好对付对手当前策略，对手就变了；其二，\emph{过拟合到训练对手}——独立 RL 学到的是对\emph{训练时遇到的特定对手}的最佳响应，换一个没见过的对手可能一败涂地。MARL 与博弈论交界的方法都在解决这两个问题——用博弈论的均衡概念（Nash/CCE）作为"鲁棒策略"的目标，用各种机制（中心化训练、种群、自博弈）处理非平稳。

\paragraph{历史：MARL 的三条主线。}
把博弈论引入 RL 最早可追溯到 Hu、Wellman 的 Nash Q-Learning\cite{hu2003nashq}——把单 agent Q-learning 的 $\max_a Q$ 换成 stage game 的 Nash，是表格式多人 RL 的奠基。之后随深度学习分化成三条主线（见 \cref{tab:sm-marl-lines}）。它们不是竞争关系，而是\emph{各管一类博弈结构}：\textbf{值分解（CTDE）管合作}（\cref{sec:sm-valuedecomp,sec:sm-policygrad} 已讲，承接 \cref{sec:sm-mpg}：合作 $=$ 团队奖励 $=$ 势函数 $=$ Markov 势博弈，故有收敛保证）；\textbf{Stackelberg RL 管层级}（leader 先动、follower 响应，用双层优化的梯度动力学，leader 的梯度要穿过 follower 的响应）；\textbf{种群博弈（PSRO）管一般和大博弈}（竞争或混合动机的大博弈，维护策略种群、不断加入"对当前种群混合策略的最佳响应"逼近均衡）。

\begin{table}[H]\centering\small
\caption{MARL 三主线全景对比}
\label{tab:sm-marl-lines}
\begin{tabular}{p{0.13\textwidth} p{0.25\textwidth} p{0.18\textwidth} p{0.13\textwidth} p{0.15\textwidth}}
\toprule
主线 & 核心机制 & 代表算法 & 均衡概念 & 适用博弈 \\
\midrule
值分解/CTDE & 中心化训练、值分解、去中心化执行 & VDN, QMIX, MAPPO & 团队最优（$=$ 势博弈纯 NE） & 合作 \\
Stackelberg RL & 双层优化梯度动力学 & Fiez 2020 等 & Stackelberg 均衡 & 层级（主从） \\
种群博弈/PSRO & 种群 $+$ 经验博弈 $+$ 最佳响应 & PSRO, JPSRO, Pipeline PSRO & Nash / CCE & 一般和（竞争/混合） \\
\bottomrule
\end{tabular}
\end{table}

\begin{note}
\textbf{Stackelberg RL 的难点：梯度要穿过对手的学习过程.}\;普通 MARL 每个 agent 假设他人固定、对当前他人策略求梯度；Stackelberg RL 里 leader 假设 follower 会\emph{重新最优响应}，所以 leader 的梯度要包含"follower 的响应如何随 leader 变化"这一项（$\partial\,\text{follower 响应}/\partial\,\text{leader 策略}$）——这需要对 follower 的优化过程求导，和逆博弈"对内层求解器求导"是同一种"对优化过程求导"的思想。\rebuilt{Gerstgrasser--Parkes（"Oracles \& Followers"）给出统一框架，并证明一个反直觉的不可能结果：当 follower 总是即时最佳响应时，leader 用 RL 学 Stackelberg 均衡可能失败，必须让 follower 有更丰富的学习动态——待核。}适用场景：机制设计（拍卖方 leader、竞拍者 follower）、人机协作（机器人 leader 影响人）、安全博弈（防御者 leader、攻击者 follower）。\pz{Stackelberg 博弈（主从均衡）的基础定义应在规划部博弈规划章给出，本节只讲其 RL 化的难点，待 \cref。}
\end{note}

\subsection{PSRO：把 Double Oracle 推广到深度 RL}
\label{sec:sm-psro-core}

要理解 PSRO，先回顾它的前身 \textbf{Double Oracle（DO）}。给定一个巨大的博弈，DO 不直接求解，而是：维护双方的\emph{策略子集}；在子集张成的\emph{小博弈}上求 Nash；对这个 Nash，各自训练一个在\emph{完整动作空间}里的最佳响应；把最佳响应加入子集；迭代。每轮要么找到能改进的最佳响应（子集扩大），要么找不到（说明当前小博弈的 Nash 已是完整博弈的 Nash）——有限博弈下必然收敛。DO 的精妙：它只在一个\emph{小}的、不断增长的策略子集上求解，避开了在完整大博弈上求 Nash。

\begin{paper}[PSRO：统一博弈论视角的多智能体强化学习]\label{paper:psro}
\textbf{问题}：独立 RL 对训练对手过拟合、泛化差；如何在大博弈上稳健地逼近均衡？
\textbf{核心思想}\cite{lanctot2017psro}：策略空间响应预言机（PSRO）是 Double Oracle 的自然推广——meta 博弈的选择对象是\emph{策略}而非\emph{动作}，且用深度 RL 训练最佳响应。流程：维护一个\emph{策略种群}（每个 agent 一组训练好的策略，初始可只含一个随机策略）；用所有策略两两对战结果构成\emph{经验博弈}（empirical game，一个收益矩阵）；在经验博弈上用 \emph{meta-solver} 求 \emph{meta 策略}（种群里策略上的混合分布 $\sigma$）；用深度 RL 训练一个新策略作为对"对手按 $\sigma$ 混合"的最佳响应（oracle）；把新策略加入种群；迭代。
\textbf{关键结果}：任何 meta-solver 都可插入；它泛化了独立 RL、迭代最佳响应、Double Oracle、虚拟自博弈等已有算法。
\textbf{与本书关系}：它是博弈与深度学习交界统一性最强的框架，承接 \cref{sec:sm-cbf-gne} 的 GNE 概念与 \cref{sec:sm-mpg} 的势博弈收敛。
\textbf{局限}：每轮训练一个 RL 策略极贵；oracle 的近似质量、收益矩阵估计方差、种群无限增长都是论文略过的工程坑。
\end{paper}

\begin{insight}[meta-solver 的选择 $=$ 一个旋钮统一所有 MARL]
PSRO 最漂亮的地方，是 meta-solver（在经验博弈上求 meta 策略的算法）可以\emph{任意替换}，而不同选择恰好对应不同的已有算法：meta-solver $=$ \emph{均匀分布}（对种群里所有策略等概率混合）$\to$ PSRO 退化成\textbf{虚拟自博弈（FSP）}；$=$ \emph{只选最后一个策略}（Dirac）$\to$ 退化成\textbf{迭代最佳响应}，进一步退化成\textbf{独立 RL}；$=$ \emph{Nash 均衡}（在经验博弈上求 Nash）$\to$ 标准 \textbf{PSRO}（即推广的 Double Oracle）；$=$ \emph{相关均衡 CCE}（$n$ 人一般和）$\to$ \textbf{JPSRO}\cite{marris2021jpsro}（用 CCE 避开 Nash 的 PPAD-hard）；$=$ \emph{$\alpha$-Rank} $\to$ $\alpha$-PSRO\cite{muller2020alpharank}。一句话：\textbf{PSRO 用"meta-solver"这一个旋钮，把独立 RL、迭代最佳响应、虚拟自博弈、Double Oracle、JPSRO 统一进一个框架}——这就是论文标题"A Unified Game-Theoretic Approach"的含义。理解了这点，你就不会把这些 MARL 算法当孤立方法，而看到它们都是 PSRO 在不同 meta-solver 下的特例。
\end{insight}

\paragraph{为什么 PSRO 收敛：可利用度。}
PSRO 继承了 Double Oracle 的收敛性，理解"为什么收敛、何时停"需要 \textbf{可利用度（exploitability）}——博弈求解里度量"离 Nash 多远"的标准工具。一个策略（或策略组合）的可利用度，定义为"对手对它做最佳响应时，它会损失多少"。Nash 均衡的可利用度为零（在 Nash 处没有任何对手的最佳响应能额外占便宜，这正是 Nash 的定义）；离 Nash 越远可利用度越高。PSRO 每一轮对当前 meta 策略训练一个\emph{最佳响应}并加入种群——这个最佳响应恰恰是"最能利用当前 meta 策略的对手"。如果它能占到便宜（可利用度 $>0$），说明当前 meta 策略还不是 Nash；如果占不到（可利用度 $\approx0$），说明已是（近似）Nash，PSRO 收敛。所以收敛判据很自然：\textbf{监控可利用度，它随迭代下降到零即收敛}。

\begin{insight}[PSRO $=$ 用"最能利用我的对手"反复打磨自己]
盯住 PSRO 的核心循环——它每轮都在找"最能利用当前策略的对手"，然后把这个对手纳入种群、让 meta 策略学会防御它。这是一种\textbf{对抗性的自我改进}：不断暴露自己的弱点（被最佳响应利用），再修补弱点（meta 策略防御），循环往复直到无懈可击（可利用度为零 $=$ Nash）。这个思想远超 PSRO 本身——它是一切"通过对抗变强"方法的共同内核：GAN（生成器 vs 判别器）、对抗训练、AlphaStar 的 League（主 agent vs 专门利用它的 exploiter）。\textbf{可利用度从正降到零的过程，就是策略从有弱点到鲁棒均衡的过程。} 种群的另一个作用是把"对单一对手过拟合"系统地转化为"对一族对手鲁棒"——新策略是对种群\emph{混合}的最佳响应，要同时应对种群里所有历史策略，于是学到的策略更鲁棒，这正是 AlphaStar 用 League 训练的核心思想。
\end{insight}

\begin{codebox}[Python]{最小 PSRO 在石头剪刀布上逼近混合 Nash}
import numpy as np
A = np.array([[0, -1, 1], [1, 0, -1], [-1, 1, 0]], float)   # rock-paper-scissors (row payoff)

def meta_nash(M):
    # meta-solver: fictitious play (FP) for mixed Nash of zero-sum subgame M
    n, m = M.shape
    rc, cc = np.zeros(n), np.zeros(m); rs, cs = np.ones(n) / n, np.ones(m) / m
    for _ in range(2000):
        rc[np.argmax(M @ cs)] += 1; rs = rc / rc.sum()         # row BR to current col mix
        cc[np.argmin(rs @ M)] += 1; cs = cc / cc.sum()         # col BR to current row mix
    return rs, cs

pop = [0]                                       # init population: only strategy 0 (rock)
for epoch in range(6):
    sub = A[np.ix_(pop, pop)]                   # empirical game: pairwise payoff submatrix
    rs, cs = meta_nash(sub)                     # meta-solver = Nash -> standard PSRO
    opp_mix = np.zeros(3)                        # map opponent meta-strategy to full action space
    for idx, a in enumerate(pop):
        opp_mix[a] += cs[idx]
    br = int(np.argmax(A @ opp_mix))            # oracle: BR to meta-strategy (real: deep RL)
    if br not in pop:
        pop.append(br)                          # add new strategy to population
# population grows to [0,1,2]; final meta-strategy ~ [0.33,0.34,0.34] (uniform Nash)
\end{codebox}

PSRO 从"只会出石头"出发：第一轮发现对手会出布克制，于是把布纳入；下一轮发现对手会出剪刀……逐轮把三个动作全纳入种群，meta 策略收敛到均匀混合——正是石头剪刀布的 Nash。这就是 PSRO 核心机制的最小兑现：\textbf{维护种群 $+$ 经验博弈求 meta 策略 $+$ 加最佳响应}，在一个独立 RL 会循环不收敛的博弈上，稳定收敛到混合 Nash。

\begin{note}
\textbf{PSRO 家族的演进（三个瓶颈 $\to$ 三个变体）.}\;\rebuilt{年份/会议待核：}标准 PSRO 用 Nash 作 meta-solver、串行训练，有三个瓶颈：(1) \emph{Nash 在 $n$ 人一般和博弈不可算}（PPAD-hard）——$\alpha$-PSRO\cite{muller2020alpharank}用 $\alpha$-Rank（基于马尔可夫链的演化排序）替代 Nash，多项式可算、适用一般和；(2) \emph{Nash 的解概念太严格}——JPSRO\cite{marris2021jpsro}用相关均衡（CCE）作 meta-solver，多项式可算、在 $n$ 人一般和里有意义，保证收敛到 CCE；(3) \emph{串行训练慢}——Pipeline PSRO\cite{mcaleer2020pipeline}用分层流水线并行训练多个最佳响应，在 Barrage Stratego（约 $10^{50}$ 状态）达到 SOTA。PSRO 家族的演进就是在"meta-solver 可算性、解概念恰当性、训练并行性"三维上逐个攻克——没有一个变体在所有维度最优，选哪个取决于博弈落在三角的哪个角。AlphaStar\cite{vinyals2019alphastar}的 League 训练是这套"种群 $+$ 对抗打磨"思想的工业级兑现，关键工程创新是给种群成员分配角色（main agent / main exploiter / League exploiter）以系统性覆盖弱点。
\end{note}

\paragraph{Nash-Q：把博弈求解嵌进 Bellman 更新。}
PSRO 在策略层面求均衡，而更早的 Nash-Q\cite{hu2003nashq}在\emph{值函数}层面引入博弈。它的核心改动一句话能说清：把单 agent Q-learning 的 $\max_a Q$ 换成 stage game 的 \textbf{Nash 值}。单 agent 假设未来 agent 会贪心地 $\max_a Q$，但多 agent 里未来是一个\emph{博弈}（所有 agent 同时选动作），不能简单 max。Nash-Q 把这个 $\max$ 换成"未来 stage game 的 Nash 值"——它假设未来所有 agent 会玩 Nash，用 Nash 的期望收益做 bootstrapping。代价是每次 Bellman 更新都要在 $Q(s',\cdot,\cdot)$ 定义的 stage game 上\emph{求一次 Nash}（贵），且只在 stage game 有唯一 Nash 时保证收敛（多 Nash 时收敛失败——又是均衡选择问题）。这限制了它只能用于小的表格式博弈，也是为什么深度时代转向 PSRO（在策略层面、用经验博弈，而非每步求 Nash）。

\begin{pitfall}[陷阱：独立 RL 直接解多智能体 / PSRO 的 meta-solver 必须用 Nash]
其一：\textbf{以为独立 RL（每个 agent 各跑各的）能直接解多智能体问题}——训练不稳、对训练对手过拟合、循环博弈里永不收敛（根本原因见 \cref{sec:sm-nonstationary}）。合作用 CTDE（QMIX/MAPPO），竞争/混合用 PSRO（种群避免过拟合单一对手）；检验方法：用没见过的对手测试，性能暴跌说明过拟合。其二：\textbf{以为 PSRO 的 meta-solver 必须用 Nash}——在 $n$ 人一般和大博弈里坚持用 Nash 会撞上 PPAD-hard 或多均衡选择问题。meta-solver 可插拔：2 人零和用 Nash（可算且有意义），$n$ 人一般和用 CCE（JPSRO 避开 Nash 的 PPAD-hard），要简单基线用 uniform（退化成 FSP）。其三：\textbf{经验博弈收益矩阵用少量对战估计}——收益是期望，局数太少噪声大致 meta 策略不稳；需足够对战局数降方差。
\end{pitfall}

\begin{practice}
\begin{enumerate}
  \item 复现石头剪刀布的 PSRO demo，扩展：(1) 换更大的对称零和博弈（如 RPS-蜥蜴-Spock，5 动作），验证收敛到混合 Nash；(2) 把 meta-solver 从 Nash 换成 uniform（退化成 FSP），对比收敛行为；(3) 把 oracle 从"精确 BR"换成"带噪声的近似 BR"（模拟深度 RL 不完美），观察对收敛的影响。
  \item 给定三个多智能体场景：(a) 10 个无人机协同搜索（共享"覆盖率"奖励）；(b) 一个拍卖机制设计（拍卖方 leader、竞拍者 follower）；(c) 6 人德州扑克。分别该用三主线中的哪条？从博弈结构（合作/层级/一般和）论证理由，并说清各自的计算代价与收敛性。
  \item 论证 PSRO 的收敛判据"可利用度降到零"为什么等价于"meta 策略是 Nash"；解释为什么独立 RL 是 PSRO 在 meta-solver $=$ Dirac（只选最后策略）下的特例。
\end{enumerate}
\end{practice}

% ════════════════════════════════════════════════════════════════
\section{统一框架、CFR 与设计空间：四类方法如何选}
\label{sec:sm-frameworks}

\cref{sec:sm-cbf-game,sec:sm-potential,sec:sm-psro} 给了博弈安全（CBF）、协调（势博弈）、学习（PSRO/值分解/策略梯度）的方法，但散落在不同实现里。研究和工程需要一个\textbf{统一框架}来定义博弈、运行算法、对比方法——这就是 OpenSpiel。本节借它讲清自博弈与 CFR 的博弈论诠释，并把全章方法放进一个统一的选型框架。

\paragraph{OpenSpiel：Game/State/Policy 三层抽象。}
OpenSpiel\cite{lanctot2019openspiel}是博弈-MARL 研究的事实标准框架（C++ 核心 $+$ Python 绑定，支持 80+ 游戏与 CFR/PSRO/AlphaZero 全家桶）。最核心的两个抽象是 \texttt{Game}（博弈定义：玩家数、动作数、收益上下界、初始状态）和 \texttt{State}（博弈状态：合法动作 \texttt{LegalActions}、执行动作 \texttt{ApplyAction}、当前玩家、是否终局、终局收益 \texttt{Returns}、信息集字符串 \texttt{InformationStateString}）。其中 \texttt{InformationStateString} 尤其关键——它是处理\emph{不完美信息博弈}（如扑克，看不到对手的牌）的核心：完美信息博弈里状态就是全部信息，但不完美信息博弈里玩家只能观测到自己的\emph{信息集}（information set，所有与自己观测一致的状态的集合），CFR 等算法正是在信息集上做后悔最小化。

\begin{insight}[统一抽象 $=$ "博弈"与"算法"解耦]
OpenSpiel 设计的精髓，是把"博弈是什么"和"用什么算法解"\textbf{彻底解耦}。\texttt{Game}/\texttt{State} 抽象只描述博弈的规则（谁动、有什么动作、终局收益），完全不涉及怎么求解；算法（CFR、PSRO、MCTS）只依赖这套抽象接口，不关心具体是哪个游戏。于是：\emph{新增一个游戏}（继承 \texttt{Game}/\texttt{State}、注册），所有算法立刻能在它上面跑；\emph{新增一个算法}（基于抽象接口实现），它立刻能跑所有 80+ 游戏。这是 $N$ 游戏 $\times$ $M$ 算法的组合爆炸被一套抽象化解——你不用写 $N\times M$ 份代码，只需 $N$ 个游戏 $+$ $M$ 个算法。这正是软件工程"面向接口编程"在博弈研究中的体现，也是 OpenSpiel 能成为事实标准的根本原因。
\end{insight}

\paragraph{CFR：no-regret 学习的平均收敛到 Nash。}
OpenSpiel 里最重要的算法之一是 \textbf{CFR（Counterfactual Regret Minimization，反事实后悔最小化）}\cite{zinkevich2008cfr}——求解不完美信息博弈（扑克）的主力，也是 Pluribus 等扑克 AI 的基础。它的核心思想：维护每个动作的\emph{累积后悔}（regret，"早知道该多出这个动作多好"），按后悔的正部比例选动作（regret matching）；理论保证\emph{平均策略}收敛到 Nash。注意一个微妙但关键的点：收敛到 Nash 的是\textbf{平均策略}，不是当前策略——当前策略（瞬时的 regret matching）会一直震荡（后悔在变），但\emph{时间平均}收敛到 Nash。

\begin{codebox}[Python]{CFR 的核心 regret matching（石头剪刀布）}
import numpy as np
A = np.array([[0, -1, 1], [1, 0, -1], [-1, 1, 0]], float)   # rock-paper-scissors (row payoff)
nA = 3
regret_sum = np.zeros((2, nA)); strategy_sum = np.zeros((2, nA))

def get_strategy(rs):                                # regret matching: proportional to positive regret
    pos = np.maximum(rs, 0); s = pos.sum()
    return pos / s if s > 0 else np.ones(nA) / nA    # all-negative regret -> uniform

for t in range(10000):
    strat = [get_strategy(regret_sum[p]) for p in range(2)]
    for p in range(2):
        strategy_sum[p] += strat[p]                  # accumulate strategy (for averaging)
    a = [np.random.choice(nA, p=strat[p]) for p in range(2)]   # sample from current strategy
    util0 = A[:, a[1]]; regret_sum[0] += util0 - util0[a[0]]   # player 0: counterfactual - actual
    util1 = -A[a[0], :]; regret_sum[1] += util1 - util1[a[1]]  # player 1 (zero-sum)
# average strategy strategy_sum/sum ~ [0.33,0.33,0.33] = RPS Nash (current osc., average converges)
\end{codebox}

\begin{insight}[均衡是涌现的，不是被直接计算的]
CFR 背后是一个深刻的博弈论定理——\textbf{在零和博弈里，如果双方都用 no-regret 算法（后悔随时间次线性增长），那么他们的平均策略收敛到 Nash}。regret matching 是一种 no-regret 算法，所以双方各自最小化后悔，平均策略就收敛到 Nash——没有谁在"求解"Nash，Nash 从双方的后悔最小化中\emph{涌现}。这与 \cref{sec:sm-potential} 势博弈的"势函数下降到极小"、\cref{sec:sm-psro} PSRO 的"种群增长逼近均衡"是同一母题的不同实现：\textbf{均衡不是被直接计算的，而是从某种简单的迭代动力学（势下降 / no-regret / 种群 BR）中涌现的}。这是它能扩展到 $10^{160}$ 信息集的扑克的根本原因。\rebuilt{CFR 家族沿三条线演进（待核年份）：CFR+（regret 截断为非负 $+$ 线性加权平均，加速一个量级）；MCCFR（蒙特卡洛 CFR，采样部分轨迹而非遍历整树）；Deep CFR（用神经网络逼近后悔，突破表格内存）。}
\end{insight}

\paragraph{exploitability 是博弈求解的客观标尺。}
求解一个博弈后怎么知道解得好不好？OpenSpiel 的标准答案是 exploitability（可利用度，\cref{sec:sm-psro-core} 已引入）。为什么用它而非"对战胜率"？对战胜率依赖\emph{拿谁当对手}——你的策略可能恰好克制测试对手（胜率高）但有其他弱点，胜率高是假象；exploitability 是\emph{对最坏情况对手}的度量——它问"最能利用你的对手能占多少便宜"，无论那个对手是谁。可利用度低意味着对\emph{任何}对手都稳健，这正是 Nash 的鲁棒性。\textbf{exploitability 把"对某对手多强"换成"对最坏对手多稳健"，是博弈求解唯一客观的进展标尺}——这是 PSRO 收敛判据、CFR 评估共用的标尺。

\paragraph{自博弈的博弈论诠释。}
OpenSpiel 把许多看似不同的"自博弈"方法统一在博弈论视角下（见 \cref{tab:sm-selfplay}）。所有自博弈方法的共同内核，是\textbf{让 agent 对弈一族对手（自己的副本/历史/种群），以逼近博弈均衡}——AlphaZero 对弈自己的当前版本、FP 对弈历史平均、CFR 双方 no-regret、PSRO 对弈种群混合。它们的差异只在"对弈谁"和"如何更新"，但都在回答同一个问题：如何不靠显式求解、而靠对弈逼近均衡。

\begin{table}[H]\centering\small
\caption{自博弈方法的博弈论诠释}
\label{tab:sm-selfplay}
\begin{tabular}{p{0.18\textwidth} p{0.50\textwidth} p{0.22\textwidth}}
\toprule
自博弈方法 & 博弈论诠释 & 适用博弈 \\
\midrule
AlphaZero & 两人零和完美信息博弈的\emph{策略迭代}（MCTS 做策略改进，自博弈数据做策略评估） & 围棋、象棋 \\
虚拟博弈（FP） & 对手用历史平均策略，自己做最佳响应；势博弈/零和收敛到 Nash & 势博弈、两人零和 \\
CFR & 双方 no-regret 学习，平均策略收敛到 Nash & 两人零和不完美信息（扑克） \\
PSRO & Double Oracle 推广：种群 $+$ 经验博弈 meta 策略 $+$ 最佳响应 & 一般和大博弈（星际） \\
\bottomrule
\end{tabular}
\end{table}

\paragraph{离散 vs 连续：博弈求解的两大分支。}
最后澄清一个本章贯穿的分工——OpenSpiel 和连续微分博弈求解器（iLQGames 类）处理\emph{两类不同的博弈}，互补而非竞争。\emph{离散扩展式博弈}（OpenSpiel）：状态和动作离散，博弈以树/信息集表示，核心挑战是巨大的状态空间（扑克 $10^{160}$）和不完美信息，方法是 CFR/PSRO/自博弈，通常离线求解；\emph{连续微分博弈}（iLQGames）：状态和动作连续，博弈以轨迹表示，核心挑战是实时性和非线性动力学，方法是迭代 LQ $+$ Riccati（与 \cref{sec:ctrl-lqr,sec:ctrl-mpc} 同源），要在线实时。\textbf{机器人规控的博弈大多是连续微分博弈，而棋牌/策略游戏的博弈是离散扩展式博弈}——选哪个框架取决于你的博弈是连续轨迹还是离散决策树。

\subsection{综合对比与选型决策}
\label{sec:sm-selection}

把全章四类方法（及它们对应的博弈类型）摆在统一维度上：

\begin{table}[H]\centering\small
\caption{四类方法的多维度对比}
\label{tab:sm-four-methods}
\begin{tabular}{p{0.14\textwidth} p{0.19\textwidth} p{0.17\textwidth} p{0.19\textwidth} p{0.19\textwidth}}
\toprule
维度 & CBF + 博弈 & 势博弈 & MARL/PSRO & OpenSpiel/CFR \\
\midrule
博弈类型 & 连续、共享安全约束（GNE） & 合作、有势结构 & 一般和、大规模 & 离散扩展式、不完美信息 \\
核心目标 & 安全（不变集） & 协调（势函数极小） & 鲁棒均衡（Nash/CCE） & Nash（信息集后悔最小） \\
求解 vs 学习 & 求解（QP）/ 学习（GCBF+） & 求解（梯度）/ 学习（MPG-PG） & 学习（深度 RL） & 求解（遍历）/ 学习（Deep CFR） \\
保证 & 硬安全（QP 可行时） & 收敛到纯 NE & 收敛到均衡（2 人零和强） & 平均策略收敛 Nash \\
可扩展性 & 手工弱 / GCBF+ 强（1024 agent） & 强（去中心化） & 弱（每轮训 RL，贵） & 中（采样/神经网络扩展） \\
实时性 & 强（QP 毫秒级） & 强 & 离线训练 & 离线求解 \\
典型应用 & swarm 避碰、安全滤波 & 覆盖、编队、任务分配 & 扑克、星际、混合博弈 & 棋牌、不完美信息博弈 \\
\bottomrule
\end{tabular}
\end{table}

\begin{insight}[一张表里藏着两条正交的轴]
全章方法的分布可以用\textbf{两条正交的轴}概括。第一条轴是\emph{博弈结构}：合作（势博弈）—共享安全约束（CBF）—一般和（MARL/PSRO）—不完美信息（CFR）。结构越特殊（合作/势），越有强保证、越可去中心化；结构越一般（一般和），越需学习、越贵。第二条轴是\emph{求解 vs 学习}：博弈小/结构清晰用求解（CBF-QP、势函数梯度、CFR 遍历），博弈大/需泛化用学习（GCBF+、MPG-PG、PSRO、Deep CFR）。一句话：\textbf{选方法 $=$ 在"博弈结构（多特殊）"和"规模（多大、是否需学习）"两条轴上定位你的问题}——结构特殊 $+$ 规模小用求解类（CBF-QP/势博弈/CFR），结构一般 $+$ 规模大用学习类（GCBF+/PSRO）。当"安全 $+$ 学习"都要时，走分层架构（上层 MARL/PSRO 学策略 $+$ 下层 CBF/GCBF+ 保安全，即 Safe MARL）。
\end{insight}

\begin{note}
\textbf{近五年前沿的定位（2021--2026）.}\;\rebuilt{年份/会议待核：}前沿沿两条轴外推。沿\emph{规模}轴：GCBF+\cite{zhang2025gcbfplus}（图 CBF $+$ GNN 推向上千 agent）、Pipeline PSRO\cite{mcaleer2020pipeline}（并行训练推向超大博弈）；沿\emph{结构一般性}轴：JPSRO\cite{marris2021jpsro}（一般和 CCE）、Markov $\alpha$-Potential Games（近似势）、MPG policy gradient 收敛理论\cite{leonardos2022mpg,ding2022mpg}；新兴的\emph{安全 $\times$ 学习}深度融合：分层 Safe MARL（上层学协作、下层 CBF 保安全）。\pz{HMARL-CBF、分布式 epigraph MARL、扩散式多机规划（MMD/SMD）等 2024–2026 工作的具体出处与年份需联网核实后再补引用；本节只点到方向，勿编造来源。}
\end{note}

\begin{pitfall}[思维陷阱：盲目追最新最复杂 / 以为"安全"和"学习"二选一]
其一：\textbf{盲目追最新、最复杂的方法}——看到 GCBF+、PSRO、Deep CFR 这些 SOTA 就默认它们一定比 CBF-QP、势博弈、CFR+ 好。用 GCBF+ 解 3 个机器人的避碰是杀鸡用牛刀（要训练、要 GNN，远不如几十行 CBF-QP）；用 PSRO 解一个明明是合作的任务（该用 QMIX）。方法的优劣是\emph{相对于问题结构}的，不是绝对的。按设计空间两条轴定位你的问题，选\emph{匹配}的方法——少 agent 安全用 CBF-QP，合作用势博弈/QMIX，大规模才上 GCBF+，一般和大博弈才上 PSRO。其二：\textbf{以为"安全"和"学习"必须二选一}——以为要么用 CBF（硬安全、不学习）、要么用 MARL（学习、无硬保证）。安全（CBF）和学习（MARL）是\emph{互补}的，可分层组合——上层 MARL 学策略（追求性能/协作），下层 CBF 滤波保安全（兜底），这正是 Safe MARL 的主流形态。检验方法：把 CBF 当 safety shield 套在 RL 策略外，既学又安全。
\end{pitfall}

\begin{practice}
\begin{enumerate}
  \item 安装 OpenSpiel，跑通示例：创建一个内置游戏（如 \texttt{kuhn\_poker}），遍历其状态树、打印合法动作/信息集/终局收益；实现一个自定义 $2\times2$ 矩阵博弈（如囚徒困境），用 CFR 求近似 Nash，用 exploitability 评估求解质量。
  \item 复现 CFR regret matching demo，扩展到 Kuhn poker（最小不完美信息扑克，3 张牌）：在信息集上做 CFR，验证平均策略收敛到已知的 Kuhn poker Nash；对比 CFR 和 CFR+ 的收敛速度（用平均策略评估，不是当前策略）。
  \item 给定五个场景：(a) 4 辆车无保护左转；(b) 200 架无人机编队穿越障碍场；(c) 6 人德州扑克；(d) 10 个仓储机器人协同搬运（共享吞吐量奖励）；(e) 一个需要既学协作又保证不碰撞的搜救 swarm。用 \cref{tab:sm-four-methods} 与本节两条轴为每个选方法，论证理由。
\end{enumerate}
\end{practice}

% ════════════════════════════════════════════════════════════════
\section*{本章常见误解汇总}
\addcontentsline{toc}{section}{本章常见误解汇总}

\begin{table}[H]\centering\small
\begin{tabular}{p{0.46\textwidth} p{0.46\textwidth}}
\toprule
\textbf{常见误解} & \textbf{正确理解} \\
\midrule
CBF 保证一定到达目标 & 只保安全（safety，不碰撞）不保活性（liveness，到达）；对称场景会死锁（\cref{sec:sm-cbf-game}）\\
CBF 安全证书不需要任何对手信息 & 免的是\emph{意图/代价}建模，但需要对手\emph{状态}（相对位置/速度）；约束 $\dot h_{ij}$ 含 $\dot x_j$ \\
多机 CBF-QP 是联立求解整个 GNE & 理论上解\emph{是} GNE（\cref{thm:sm-cbf-gne}），实现上是去中心化近似（$u_j$ 当已知代入）\\
任何多机协调问题都是势博弈 & 势博弈要存在全局势函数（强条件）；合作是充分条件之一，需\emph{验证}或主动设计 \\
收敛到纯 Nash 就是全局最优 & 纯 Nash $=$ 势函数\emph{局部}极小，非全局；需多次初始化/退火跳出（\cref{thm:sm-potential-props}）\\
MARL 就是把单智能体 RL 复制 $N$ 份 & 那是独立学习，死于非平稳性；MARL 是"联合优化、分布执行"的协调问题 \\
全局状态 $s$ $=$ 观测拼接 $(o_1,\dots,o_N)$ & 通常\emph{不等}（观测部分、可能重叠/有盲区）；仅特殊情况相等 \\
非平稳性可靠调小学习率解决 & 框架性问题（追逐移动靶），调小步长只是移动靶和追逐者一起慢；须改信息结构（CTDE）\\
找到纳什均衡就是找到最优 & Nash 只保"无人单方面变好"，不保集体最优；囚徒困境是经典反例 \\
有了 CTDE 就自动解决信用分配 & CTDE 治非平稳/维度爆炸；信用分配正交，靠反事实基线/值分解 \\
个体 utility $Q_i$ 是智能体 $i$ 的真实价值 & $Q_i$ 是 utility 非 $Q$ 值，无独立语义、尺度可重标定；只用于 $\arg\max_{a_i}Q_i$ \\
QMIX 表达力够强能拟合任意协作任务 & 单调性约束表达不了"非单调耦合"（动作好坏随他人翻转），需 QTRAN/QPLEX \\
MAPPO 的 critic 吃联合动作 $Q(s,\mathbf{a})$ & MAPPO 的 critic 是状态值 $V_\phi(s)$（不吃动作），优势用 GAE（\cref{eq:sm-mappo-gae}）\\
HAPPO 顺序更新只是"循环更新每个 agent" & 必须每更新完一个 agent 就重算后续的重要性比，否则失去单调改进保证 \\
独立 RL（各跑各的）能解多智能体 & 非平稳 $+$ 过拟合训练对手；合作用 CTDE、竞争用 PSRO \\
PSRO 的 meta-solver 必须用 Nash & meta-solver 可插拔；$n$ 人一般和用 CCE（JPSRO 避 PPAD-hard）、uniform $\to$ FSP \\
CFR 的当前策略收敛到 Nash & 收敛的是\emph{时间平均}策略；当前策略持续震荡，评估/使用平均策略 \\
OpenSpiel 能解机器人连续控制博弈 & OpenSpiel 面向离散扩展式博弈（棋牌）；连续微分博弈用 iLQGames（\cref{sec:ctrl-mpc} 同源）\\
追最新最复杂方法一定更好 & 方法优劣\emph{相对于问题结构}；少 agent/合作/小规模用简单法更合适 \\
"安全"和"学习"必须二选一 & 互补，可分层（上层 MARL 学 $+$ 下层 CBF 滤波保安全 $=$ Safe MARL）\\
\bottomrule
\end{tabular}
\end{table}

% ════════════════════════════════════════════════════════════════
\section*{本章小结}
\addcontentsline{toc}{section}{本章小结}

本章用"博弈作为统一语言"把三个看似分立的领域串成一体：\textbf{安全控制、多机协调、多智能体学习，本质都是博弈，只是博弈的类型不同}。\cref{sec:sm-cbf-game} 把"安全证书"表述成广义 Nash 均衡——多机最小范数 CBF-QP（\cref{eq:sm-cbf-qp}）的 KKT 条件联合等价于 GNE（\cref{thm:sm-cbf-gne}），GCBF+ 用图神经网络把证书学出来、靠"成对碰撞的局部性"扩展到上千 agent。\cref{sec:sm-potential} 把"分布式协调"表述成博弈——势函数（\cref{eq:sm-potential-def}）使自私优化等价于集体降同一全局目标，纯 Nash 存在且分布式学习收敛（\cref{thm:sm-potential-props}），Voronoi 覆盖是典范，Markov 势博弈则为合作 MARL 的收敛性奠基。\cref{sec:sm-marl-form,sec:sm-nonstationary,sec:sm-ctde} 给出 MARL 的公共语言（Markov Game/Dec-POMDP）、第一性困难（非平稳性破坏压缩映射）与破局范式（CTDE 把漂移从环境挪进条件，\cref{eq:sm-central-bellman}）。\cref{sec:sm-valuedecomp,sec:sm-policygrad} 是两条硬核技术路线——值分解的 IGM 原则与单调混合（\cref{thm:sm-mono-igm}）、策略梯度的反事实基线与多智能体优势分解引理（\cref{thm:sm-advantage-decomp}）。\cref{sec:sm-psro} 把"多智能体学习"表述成求博弈均衡，PSRO 用可插拔的 meta-solver 统一诸方法。\cref{sec:sm-frameworks} 给出操作这些博弈的统一框架与"结构 $\times$ 规模"两条轴的选型。一条更深的主线：\textbf{均衡常常不是被直接计算的，而是从简单的迭代动力学中涌现的}——势函数单调下降、PSRO 种群增长、CFR 双方 no-regret 的平均，都是同一哲学的不同实现。

\begin{table}[H]\centering\small
\caption{符号表}
\begin{tabular}{lll}
\toprule
符号 & 含义 & 首次出现 \\
\midrule
$h(x),\ \mathcal{C}=\{h\ge0\}$ & CBF 标量函数 / 安全集 & \cref{sec:sm-cbf-game} \\
$h_{ij}=\|p_i-p_j\|^2-D_s^2$ & 机间安全裕度 & \cref{eq:sm-cbf-qp} \\
$\alpha$ & 扩展 class-$\mathcal{K}$ 函数（线性）系数 & \cref{eq:sm-cbf-cond} \\
$\lambda_{ij}\ge0$ & CBF-QP 的 KKT 乘子 & \cref{thm:sm-cbf-gne} \\
$\Phi$ & 势函数（势博弈/MPG） & \cref{eq:sm-potential-def} \\
$V_i$（Voronoi 格） & 机器人 $i$ 的 Voronoi 格 & \cref{eq:sm-voronoi} \\
$N,\ i,\ -i$ & 智能体数 / 下标 / 除 $i$ 外 & \cref{def:sm-markovgame} \\
$\mathbf{a}=(a_1,\dots,a_N)$ & 联合动作 & \cref{def:sm-markovgame} \\
$\boldsymbol{\pi},\ \pi_i(a_i\mid o_i)$ & 联合策略 / 分布式策略 & \cref{def:sm-markovgame} \\
$\tau_i$ & 局部动作-观测历史 & \cref{def:sm-decpomdp} \\
$P_i^{\boldsymbol{\pi}_{-i}}$ & 智能体 $i$ 的有效转移核 & \cref{eq:sm-effective-core} \\
$Q_i(o_i,a_i)$ & 个体 utility（值分解） & \cref{eq:sm-vdn} \\
$Q_{\text{tot}}(s,\mathbf{a})$ & 联合动作值 & \cref{eq:sm-vdn} \\
$V_\phi(s)$ & 集中式状态值函数（MAPPO） & \cref{eq:sm-mappo-gae} \\
$A_i,\ A^{i_{1:k}}$ & 个体优势 / 条件优势 & \cref{eq:sm-coma,eq:sm-adv-decomp} \\
$\rho_i$ & 重要性采样比 & \cref{eq:sm-mappo-clip} \\
$D_i$ & 差分奖励 & \cref{eq:sm-diff-reward} \\
$\phi_i$ & Shapley 值 & \cref{eq:sm-shapley} \\
$\sigma$ & 经验博弈上的 meta 策略 & \cref{paper:psro} \\
$\gamma,\ \eta,\ \lambda$ & 折扣因子 / 学习率 / GAE 参数 & 全章 \\
\bottomrule
\end{tabular}
\end{table}

\begin{table}[H]\centering\small
\caption{定理与关键式速查表}
\begin{tabular}{lll}
\toprule
定理 / 公式 & 一句话说明 & 对应 \\
\midrule
多机 CBF-QP \cref{eq:sm-cbf-qp} & 最小范数安全修正，约束含邻居运动 $\dot x_j$ & \cref{sec:sm-cbf-gne} \\
CBF-QP $=$ GNE \cref{thm:sm-cbf-gne} & 共享约束使联合 KKT 即 GNE 的 KKT & \cref{sec:sm-cbf-gne} \\
relaxed CBF \cref{eq:sm-relaxed-cbf} & 软约束 $+$ 松弛保 QP 永远可行 & \cref{sec:sm-cbf-code} \\
势函数定义 \cref{eq:sm-potential-def} & 单边偏离收益变化 $=$ $\Phi$ 变化 & \cref{def:sm-potential} \\
势博弈三性质 \cref{thm:sm-potential-props} & 纯 NE 存在 $+$ 更优响应收敛 & \cref{sec:sm-potential} \\
Voronoi locational cost \cref{eq:sm-voronoi} & Lloyd $=$ 势函数坐标下降 & \cref{sec:sm-potential} \\
有效转移核 \cref{eq:sm-effective-core} & 非平稳性的精确来源 & \cref{def:sm-nonstationarity} \\
集中条件贝尔曼 \cref{eq:sm-central-bellman} & 以联合动作为条件恢复平稳性 & \cref{sec:sm-ctde} \\
IGM 原则 \cref{eq:sm-igm} & 局部 argmax $=$ 全局 argmax & \cref{def:sm-igm} \\
单调 $\Rightarrow$ IGM \cref{thm:sm-mono-igm} & QMIX 非负权重的理论根据 & \cref{sec:sm-valuedecomp} \\
MADDPG 梯度 \cref{eq:sm-maddpg} & actor 局部、critic 全局 & \cref{sec:sm-policygrad} \\
COMA 反事实基线 \cref{eq:sm-coma} & 固定别人、平均掉自己 $\to$ 信用分配 & \cref{sec:sm-policygrad} \\
MAPPO clip \cref{eq:sm-mappo-clip} & PPO 逐字搬到多智能体 $+$ 全局 $V$ & \cref{sec:sm-policygrad} \\
优势分解引理 \cref{thm:sm-advantage-decomp} & 联合优势 $=$ 逐个条件优势之和 & \cref{sec:sm-policygrad} \\
差分奖励 \cref{eq:sm-diff-reward} & 贡献 $=$ 反事实差值（金标准） & \cref{sec:sm-credit} \\
Shapley 值 \cref{eq:sm-shapley} & 归因的理论标尺（全顺序平均） & \cref{sec:sm-credit} \\
PSRO meta-solver & 一个旋钮统一 InRL/FSP/DO/JPSRO & \cref{sec:sm-psro-core} \\
CFR / no-regret & 双方后悔最小化的平均策略收敛 Nash & \cref{sec:sm-frameworks} \\
\bottomrule
\end{tabular}
\end{table}

\begin{table}[H]\centering\small
\caption{知识点总表（难度用 \texorpdfstring{$\bigstar$}{star}）}
\begin{tabular}{clll}
\toprule
\# & 知识点 & 核心要点 & 难度 \\
\midrule
1 & 两种安全 & 软（基于预测）vs 硬（基于不变集，不依赖对手模型） & \(\bigstar\bigstar\) \\
2 & CBF-QP $\leftrightarrow$ GNE & 多机最小范数 QP 的 KKT 联合 $=$ 广义 Nash 均衡 & \(\bigstar\bigstar\bigstar\) \\
3 & GCBF+ & GNN 学图 CBF；成对碰撞局部性 $\Rightarrow$ 单一证书任意规模 & \(\bigstar\bigstar\bigstar\) \\
4 & 死锁与 relaxed CBF & CBF 只保安全会死锁；软约束保 QP 可行 & \(\bigstar\bigstar\) \\
5 & 势博弈与势函数 & 自私优化 $=$ 集体降 $\Phi$；纯 NE 存在 $+$ 收敛 & \(\bigstar\bigstar\bigstar\) \\
6 & Voronoi 覆盖 & locational cost 为势函数，Lloyd $=$ 坐标下降 & \(\bigstar\bigstar\) \\
7 & Markov 势博弈 & 合作 MARL $=$ 势博弈，policy gradient 收敛 & \(\bigstar\bigstar\bigstar\) \\
8 & Markov Game / Dec-POMDP & 两套形式化；合作/竞争/混合三类 & \(\bigstar\bigstar\bigstar\) \\
9 & 非平稳性 & 移动靶破坏压缩映射；调学习率治不了 & \(\bigstar\bigstar\bigstar\bigstar\) \\
10 & CTDE & 把漂移从环境挪进条件；训练开天眼、执行接地气 & \(\bigstar\bigstar\bigstar\) \\
11 & IGM 与 VDN/QMIX & 局部 argmax $=$ 全局 argmax；单调蕴含 IGM & \(\bigstar\bigstar\bigstar\bigstar\) \\
12 & QTRAN/QPLEX & 放松单调以逼近 IGM 充要刻画 & \(\bigstar\bigstar\bigstar\bigstar\) \\
13 & MADDPG/COMA & 集中 critic 治非平稳；反事实基线解信用 & \(\bigstar\bigstar\bigstar\bigstar\) \\
14 & MAPPO & 共享 PPO $+$ 全局 $V$，最强基线 & \(\bigstar\bigstar\bigstar\) \\
15 & HAPPO 优势分解引理 & 联合优势 $=$ 逐个条件优势；顺序更新保单调改进 & \(\bigstar\bigstar\bigstar\bigstar\bigstar\) \\
16 & 信用分配三解法 & 差分奖励/反事实基线/值分解 $=$ 反事实归因 & \(\bigstar\bigstar\bigstar\) \\
17 & Shapley 值 & 归因的理论标尺；指数复杂度需近似 & \(\bigstar\bigstar\bigstar\bigstar\) \\
18 & MARL 三主线 & 值分解（合作）/Stackelberg（层级）/PSRO（一般和） & \(\bigstar\bigstar\bigstar\) \\
19 & PSRO & Double Oracle 推广；meta-solver 可插拔统一诸方法 & \(\bigstar\bigstar\bigstar\bigstar\) \\
20 & CFR 与 exploitability & no-regret 平均收敛 Nash；可利用度是客观标尺 & \(\bigstar\bigstar\bigstar\) \\
21 & 设计空间两条轴 & 博弈结构（多特殊）$\times$ 规模（求解 vs 学习）定选型 & \(\bigstar\bigstar\bigstar\) \\
\bottomrule
\end{tabular}
\end{table}

% ════════════════════════════════════════════════════════════════
\section*{延伸阅读}
\addcontentsline{toc}{section}{延伸阅读}
\begin{itemize}
  \item \textbf{多机安全证书}：Wang--Ames--Egerstedt\cite{wang2017safety}（多机安全屏障证书的奠基，本章 \cref{sec:sm-cbf-game} 的源）；Ames 等\cite{ames2017cbf}（CBF 与控制 Lyapunov 函数的统一框架）；GCBF+\cite{zhang2025gcbfplus}（神经图 CBF，可扩展安全的代表）。
  \item \textbf{势博弈与合作控制}：Monderer--Shapley\cite{monderer1996potential}（势博弈奠基）；Marden--Arslan--Shamma\cite{marden2009cooperative}（势博弈用于合作控制）；Leonardos 等\cite{leonardos2022mpg}、Ding 等\cite{ding2022mpg}（Markov 势博弈的策略梯度收敛理论）。
  \item \textbf{MARL 形式化与综述}：Littman\cite{littman1994markov}（Markov Game 形式化奠基）；Bernstein 等\cite{bernstein2002decpomdp}（Dec-POMDP 与 NEXP-完全复杂度）。
  \item \textbf{值分解}：Sunehag 等\cite{sunehag2018vdn}（VDN）；Rashid 等\cite{rashid2018qmix}（QMIX 单调混合）；Son 等\cite{son2019qtran}（QTRAN）；Wang 等\cite{wang2021qplex}（QPLEX 的 IGM 充要刻画）。
  \item \textbf{多智能体策略梯度}：Lowe 等\cite{lowe2017maddpg}（MADDPG）；Foerster 等\cite{foerster2018coma}（COMA 反事实基线）；Yu 等\cite{yu2022mappo}（MAPPO 强基线翻案）；Kuba 等\cite{kuba2022happo}（HAPPO 优势分解引理）；Zhong 等\cite{zhong2024harl}（HARL 异构库）。
  \item \textbf{种群博弈与博弈求解框架}：Lanctot 等\cite{lanctot2017psro}（PSRO）；Marris 等\cite{marris2021jpsro}（JPSRO/CCE）；McAleer 等\cite{mcaleer2020pipeline}（Pipeline PSRO）；Hu--Wellman\cite{hu2003nashq}（Nash-Q）；Zinkevich 等\cite{zinkevich2008cfr}（CFR）；Lanctot 等\cite{lanctot2019openspiel}（OpenSpiel 框架）；Vinyals 等\cite{vinyals2019alphastar}（AlphaStar，PSRO 思想的工业级兑现）。
\end{itemize}

\section*{本章与后续章节的关系}
\addcontentsline{toc}{section}{本章与后续章节的关系}
\begin{table}[H]\centering\small
\begin{tabular}{lll}
\toprule
章节 & 关系 & 本章铺垫 / 复用 \\
\midrule
\cref{ch:planning} 运动规划导论 & 规划器给标称控制，本章 CBF 层做安全滤波 & 搜索/采样/Bellman（\cref{sec:plan-dp}）\\
\cref{sec:ctrl-lyapunov} Lyapunov 稳定性 & CBF/势函数收敛的同源逻辑 & 势函数 $=$ Lyapunov 函数、单调下降必收敛 \\
\cref{sec:ctrl-lqr,sec:ctrl-mpc} LQR/MPC & 连续微分博弈求解 $\leftrightarrow$ 离散扩展式博弈 & iLQ/Riccati 与 OpenSpiel 的分工（\cref{sec:sm-frameworks}）\\
控制部 安全控制章 & CBF/CLF、HJ 可达性、鲁棒/随机 MPC 深讲 & 本章只给规划用途，深推待 \pz 标注的控制部 \\
规划部 博弈规划章 & 逆博弈、iLQGames、Stackelberg 均衡基础 & 本章承接其安全兜底与 RL 化（\cref{sec:sm-cbf-game,sec:sm-psro}）\\
\bottomrule
\end{tabular}
\end{table}

\section*{故障排查手册}
\addcontentsline{toc}{section}{故障排查手册}
\begin{enumerate}
  \item \textbf{症状}：多机 CBF-QP 在对称场景\emph{僵住不动到不了目标}。\textbf{原因}：死锁（CBF 只保 safety 不保 liveness，\cref{sec:sm-cbf-game}）。\textbf{对策}：对称破缺（切向偏置/随机扰动）、优先级、或高层协调。
  \item \textbf{症状}：密集场景 CBF-QP \emph{无可行解/返回 NaN}。\textbf{原因}：控制有界下多约束冲突、可行域为空。\textbf{对策}：用 relaxed CBF（\cref{eq:sm-relaxed-cbf}），监控松弛量、过大时紧急制动。
  \item \textbf{症状}：独立学习\emph{训练曲线剧烈震荡、不收敛}，换学习率/网络都没用。\textbf{原因}：非平稳性（移动靶，\cref{sec:sm-nonstationary}）。\textbf{对策}：改信息结构——合作用 CTDE（集中 critic 喂全局状态/联合动作），不是调超参。
  \item \textbf{症状}：多智能体 off-policy 训练\emph{越训越差、增大 buffer 反而恶化}。\textbf{原因}：replay 数据过期（旧经验来自他人旧策略）。\textbf{对策}：把联合动作存进 buffer 并作集中 critic 输入（MADDPG），或改用 on-policy 的 MAPPO（不存 buffer）。
  \item \textbf{症状}：QMIX \emph{训练 loss 在降但执行时智能体协调不起来}。\textbf{原因}：混合网络权重未加非负约束 $\Rightarrow$ IGM 被破坏。\textbf{对策}：超网络输出权重必须过 \texttt{torch.abs}/\texttt{F.relu}（偏置不需要），单元测试数值检查 $\partial Q_{\text{tot}}/\partial Q_i\ge0$。
  \item \textbf{症状}：QMIX 在某任务上\emph{系统性把最优联合动作学偏}，调参无效。\textbf{原因}：任务有非单调耦合，超出单调性表达力。\textbf{对策}：换 QTRAN/QPLEX 或策略梯度（MAPPO 不受 IGM 约束）；先用小矩阵博弈测方法的表达力上限。
  \item \textbf{症状}：参数共享的 MAPPO 智能体\emph{无法分化角色}（需领航-跟随却行为相同）。\textbf{原因}：参数共享 $+$ 对称观测 $\Rightarrow$ 对称策略。\textbf{对策}：在 actor 输入拼接 agent one-hot ID 或可学习 embedding 打破对称。
  \item \textbf{症状}：HAPPO \emph{不如 MAPPO 稳、体现不出异构优势}。\textbf{原因}：顺序更新未重算重要性比（退化成简单循环）。\textbf{对策}：每更新完一个 agent，后续 agent 的重要性比/优势按已更新前序重算（HARL 的 factor 累积逻辑），且用随机顺序。
  \item \textbf{症状}：某个本该有用的智能体\emph{退化成"什么都不做"}。\textbf{原因}：懒惰智能体（高方差被团队学会"按住"）。\textbf{对策}：用能正确归因边际贡献的信用分配（反事实/值分解），保证充分探索，监控行为熵。
  \item \textbf{症状}：PSRO 在 $n$ 人一般和博弈\emph{跑不动/不收敛}。\textbf{原因}：Nash meta-solver 撞 PPAD-hard 或多均衡选择。\textbf{对策}：换 CCE meta-solver（JPSRO）或 $\alpha$-Rank；确保经验博弈收益矩阵用足够对战局数估计。
  \item \textbf{症状}：跑 CFR 时\emph{当前策略一直震荡，以为没收敛}。\textbf{原因}：CFR 收敛的是\emph{时间平均}策略，不是当前策略。\textbf{对策}：评估/使用平均策略（\texttt{strategy\_sum} 归一化），监控其 exploitability 是否随迭代下降。
\end{enumerate}
